%% 第26回編集委員会（物理第4巻）：組 電W8
%% p4-14.tex  電-14　4つの式にまとめる：マクスウェル方程式（★3〜5・3段）
%% 中身は計画書 tools/phys4-plan.md の B「電-14」。公式の持ち主は A（A83 変位電流とマクスウェル方程式、A84 ポインティングベクトル）。
%% 数値はすべて Python で計算した：checks/物理第4巻/電-14/calc_p4_14.py（出力 calc_p4_14_out.txt）、
%%   fdtd_p4_14.py（⑩の表とプログラムの出力。出力 fdtd_p4_14_out.txt）。
%% 記号（この章。巻全体は計画書 C-1・第26回 Y-2〜Y-4）：
%%   第4部なので面積は A（ポインティングベクトル S とぶつけない：Y-4）。閉じた面・面のタイルも dA と書く（8-11・8-14 の dS と同じもの）。
%%   電場 E、磁場（磁束密度）B、磁場の強さ H（電-7 ③）、電束密度 D（電-3 ⑧）、電荷 Q（袋の中は Q_中）・q（動く電荷）、電荷密度 ρ（8-14 と同じ。
%%   この章には抵抗率が出ないので、電-4 の ρ_q とは書かない）、電流 I（瞬時も I。回路の章ではないので i は使わない）、電流密度 j、
%%   電場の流束 Φ_E・磁束 Φ_B（8-14 と同じ）、変位電流 I_d、ポインティングベクトル S、エネルギー密度 u、袋の中の電磁場のエネルギー U_中、
%%   電圧 V（体積の積分は書かず、体積の V とぶつけない）、極板の半径 R・すき間 d、道の半径 r、導線の半径 a、同軸の外の導体の内側の半径 b、
%%   長さ l、導電率 σ、比誘電率 ε_r、角周波数 ω（例題3だけ）、周波数 f、速さ v（⑨の動く人・棒）、光速 c（第4部）、
%%   ⑩の FDTD のマスの大きさ Δx・時間の1歩 Δt、波のインピーダンス η0（⑩で h＝η0H と書くためだけに使う。くわしくは電-15 ⑦）。
%% 公式の持ち主：ガウスの法則は電-2 ③（微分形 ⑧）、div B＝0・アンペールの法則は電-6 ⑥⑦、ファラデーの法則は電-8 ③（微分形 ⑨）、
%%   D は電-3 ⑧、H は電-7 ③、電気のエネルギー密度は電-3 ⑦、磁気のエネルギー密度は電-8 ⑧、P＝VI は電-4 ⑤、連続の式は電-4 ⑥、
%%   動く導体棒の起電力は電-8 ④、ローレンツ力は電-6 ③、C＝ε0S/d は電-3 ④（この章では面積を A と書く）。どれも導き直さず指す。
%%   計算が済んでいる所（8-14 ①③④⑤、8-11 ⑦⑬・S6、8-13 ⑧〜⑩、8-12 ⑧）は数学の章を指す（Y-9）。
%%   この章で導く中学・高校の公式はない（A83・A84 は大学の初め）。K-8 のとおり「この章で導く中学・高校の公式」の欄は置かない。
%% 前の巻との約束（計画書 A-6）：P13（第1巻の前付の地図のエネルギーの行「ポインティングベクトル → 電-14」）は⑦⑧、
%%   P22（第2巻の前付）・P44（第3巻の前付）も⑦。数学 8-14 の「電-14・15 の予告」は①⑤で受けた。
%% 同じ題材の数（計画書 C-4 の2）：ドライヤー 100 V・12 A・1200 W（例題5）、同軸 0.45 mm・1.6 mm・ε_r 2.3・50 Ω（例題6。半径として使う）、
%%   電-8 の棒 2 m/s・0.1 T・0.5 m・0.1 V（例題8）、8-5⑥ の 1 nC・0.1 m・899 N/C と 8-10⑪ の 10 A・5 cm・4×10⁻⁵ T（例題9）。
%% 出典（第26回 Y-6 2章の16。11月5日に確かめた）：
%%   ・海水の導電率：塩分35・15 ℃・0 気圧の標準海水で 42.914 mS/cm＝4.29 S/m（TEOS-10 の gsw_C3515 の説明 teos-10.org、
%%     Culkin・Smith 1980。検索で2つの資料が一致）。本文は「4.3 S/m とします」。海水の比誘電率は水の値 80（IAPWS の式で 20 ℃ 80.2：
%%     計画書 D17）を「とします」で使う（海水は少し小さい）。
%%   ・マクスウェル 1865年 Phil. Trans. 155 "A dynamical theory of the electromagnetic field"（計画書 D24：DOI 10.1098/rstl.1865.0008 の登録先の題）。
%%   ・ポインティング 1884年 Phil. Trans. 175 "On the transfer of energy in the electromagnetic field"（DOI 10.1098/rstl.1884.0016 の登録先の題で確かめた）。
%%   ・ディラックの磁気単極子 1931年 Proc. R. Soc. A 133 "Quantised singularities in the electromagnetic field"（DOI 10.1098/rspa.1931.0130 の登録先の題）。
%%   ・磁気単極子の探索：MoEDAL 実験、LHC の鉛どうしの衝突（2018年11月）で探し、見つからなかった
%%     （arXiv:2106.11933、Nature 602, 63 (2022)、DOI 10.1038/s41586-021-04298-1。arxiv.org で要旨を読んだ）。
%%   ・FDTD 法：K. S. Yee、IEEE Trans. Antennas Propag. 14 (1966)（DOI 10.1109/TAP.1966.1138693 の登録だけ確かめた。本文には年だけ書く）。
%%   ・ヘルツの実験 1888年ごろ：数学 8-14 のコラムと同じ書き方にそろえた。
%% 漢字：常用漢字だけ。checks/物理第4巻/電-14/kanji_check.py で確かめた。
\makeatletter
% 円い極板のコンデンサー（横から見た図）  \sxocap{(x,y)}{高さの半分}{すき間の半分}
\newcommand\sxocap[3]{\begin{scope}[shift={#1}]
  \fill[bkGray!70] ({-#3-1.6},{-#2}) rectangle ({-#3},{#2});
  \fill[bkGray!70] ({#3},{-#2}) rectangle ({#3+1.6},{#2});
\end{scope}}
% 乾電池（横向き。左が＋）  \sxobat{(x,y)}
\newcommand\sxobat[1]{\begin{scope}[shift={#1}]
  \filldraw[fill=bkOrange!25, draw=bkInk, line width=0.6pt, rounded corners=0.6mm] (-7,-3) rectangle (7,3);
  \filldraw[fill=bkGray!40, draw=bkInk, line width=0.6pt] (-8.4,-1.2) rectangle (-7,1.2);
  \node[font=\sffamily\footnotesize] at (-4.6,0) {＋};\node[font=\sffamily\footnotesize] at (4.6,0) {－};
\end{scope}}
\makeatother
\tikzset{sxo lab/.style={font=\sffamily\footnotesize, inner sep=1pt},
  sxo E/.style={-{Stealth[length=2.6pt, width=2.4pt]}, line width=0.9pt, bkRed!80!black},
  sxo B/.style={-{Stealth[length=2.6pt, width=2.4pt]}, line width=0.9pt, bkBlue!80!black, dash pattern=on 2.4pt off 1.4pt},
  sxo S/.style={-{Stealth[length=3.2pt, width=3pt]}, line width=1.4pt, bkOrange!85!black},
  sxo I/.style={-{Stealth[length=3pt, width=2.8pt]}, line width=1.1pt, bkInk},
  sxo dim/.style={{Stealth[length=2.4pt]}-{Stealth[length=2.4pt]}, line width=0.5pt, bkGray!80!black}}

\phchapter{14}{4つの式にまとめる：\\マクスウェル方程式}{電気と磁気は、別々のものなの？}{\bklevelrange{3}{5}}{4つの式にまとめる：マクスウェル方程式}
\bkmichishirube{8-8,8-9,8-10,8-11,8-13,8-14,電-2,電-6,電-8}{電-15,原-11}

\begin{goal}
\begin{itemize}[label={\color{bkTheme}●}, leftmargin=1.6\zw, itemsep=0pt, topsep=1pt]
\item \bklevel{3}\enspace 充電しているコンデンサーの極板の間を流れる変位電流と、その周りの磁場を計算できる。
\item \bklevel{4}\enspace マクスウェル方程式を積分形と微分形で読み、電磁場のエネルギーの流れ（ポインティングベクトル）を計算できる。
\item \bklevel{5}\enspace 動く人から見ると電場と磁場が入れかわることを説明し、電磁場を数値で解く方法（FDTD）を表計算で動かせる。
\end{itemize}
\end{goal}
%% 整合（K-8）：この章で導く中学・高校の公式はないので、「この章で導く中学・高校の公式」の欄は置かない。

\begin{hatena}
\noindent\begin{minipage}[c]{0.58\linewidth}\raggedright
磁石は、くぎやクリップを引きつけます。下じきでこすった髪の毛は、下じきに引きつけられます。どちらも、離れていてもはたらく力です。でも、磁石は紙切れを引きつけず、下じきはくぎを引きつけません。電気の力と磁気の力は、まったく別のものなのでしょうか。それとも、どこかでつながっているのでしょうか。
\end{minipage}\hfill
\begin{minipage}[c]{0.39\linewidth}\centering
\begin{tikzpicture}[x=1mm, y=1mm]
  % 磁石とクリップ
  \fill[bkRed!75!black] (0,20) rectangle (10,26); \fill[bkBlue!70!black] (10,20) rectangle (20,26);
  \node[text=white, font=\sffamily\bfseries\footnotesize] at (5,23) {N};\node[text=white, font=\sffamily\bfseries\footnotesize] at (15,23) {S};
  \draw[bkGray!80!black, line width=0.8pt, rounded corners=0.6mm] (23,21) rectangle (29,25) (24,22) rectangle (28.5,24);
  % 下じきと髪
  \filldraw[fill=bkBlue!12, draw=bkInk, line width=0.6pt] (0,2) rectangle (20,12);
  \node[sxo lab] at (10,7) {下じき};
  \foreach \x in {24,26,28,30} {\draw[bkInk, line width=0.5pt] (\x,0) .. controls (\x-1,5) .. (\x-3.5,9);}
  \node[sxo lab, anchor=north] at (27,-0.6) {髪の毛};
  \node[font=\sffamily\bfseries\Large, text=bkTheme] at (40,14) {？};
\end{tikzpicture}
\end{minipage}
\end{hatena}

\noindent 第1部と第2部では、電気と磁気の法則を1つずつ学んできました。電荷の周りの電場はガウスの法則（電-2 ③）、電流の周りの磁場はアンペールの法則（電-6 ⑥）、磁場が変わると生まれる電場はファラデーの法則（電-8 ③）で決まりました。この章では、これらを4つの式に並べてみます。すると、1か所だけ、つじつまの合わない所が見つかります。そこを直すと、電気と磁気は1つにまとまり、「光（電磁波）」が生まれる準備がととのいます（次の電-15）。

%% ---------------- やってみよう・思い出そう ----------------
\begin{yattemiyou}[エネルギーは、どこを通って豆電球へ届く？]
\textbf{用意するもの}：乾電池1個、豆電球（ソケットつき）、導線2本（\qty{30}{cm}くらい）、方位磁針、紙と色えんぴつ
\begin{enumerate}[label={\bkcircled[bkGreen]{\arabic*}}, leftmargin=2.4\zw, itemsep=1pt, topsep=2pt]
\item 乾電池と豆電球を導線でつなぎ、豆電球をつける。導線は、机の上で大きな輪になるように広げる。
\item 方位磁針を、導線のすぐ上や下に置いて、針の向きを調べる（電流の周りの磁場：電-6）。電流が流れる向きを逆にすると、針はどうなるかな？
\item 紙に回路をかき、導線の周りの磁場の向き（輪の矢印）と、乾電池の＋側の導線から－側の導線へ向かう電場の向き（まっすぐな矢印）をかきこむ。
\item ③の2つの矢印を、右ねじで「電場から磁場へ」回すと、ねじはどちらへ進むかな？　その向きを、ほかの色の矢印でかきこもう。
\end{enumerate}
\end{yattemiyou}

\noindent ④でかいた矢印は、乾電池から豆電球のほうを向いているはずです。電気のエネルギーは、導線の\textbf{中}ではなく、導線の周りの\textbf{空間}を通って、豆電球へ届いているのです。そのわけは⑦⑧で説明します。

\begin{omoidasou}
\bkrecall{電-2}{ガウスの法則：閉じた袋から出ていく電場の流束 $\Phi_E$ は、袋の中の電荷を $\varepsilon_0$ でわったもの（電-2 ③）。}
\bkrecall{電-6}{アンペールの法則：電流を囲む道にそって $\vec B$ をたすと、$\mu_0\times$（道を通りぬける電流）になる（電-6 ⑥）。まっすぐな電流から $r$ の所で $B=\mu_0I/(2\pi r)$（電-6 ⑤）。}
\bkrecall{電-8}{ファラデーの法則：道を1周する起電力は、道を通る磁束の減る速さ（電-8 ③）。}
\bkrecall{電-3}{平行板コンデンサー：$Q=CV$、$C=\varepsilon_0\times(\text{面積})/d$、極板の間は一様な電場 $E=V/d$（電-3 ④）。電場のエネルギー密度は $\tfrac12\varepsilon_0E^2$（電-3 ⑦）。}
\bkrecall{8-14}{4つの式の表（8-14①）。積分形と微分形は、ガウスの定理・ストークスの定理で行き来できる（8-14③）。}
\end{omoidasou}

\noindent この章は3段に分かれています。①〜④（目でわかる・手でわかる）は、コンデンサーの式 $Q=CV$ と、これまでの章の法則だけで読めます。⑤〜⑩（式でわかる）では、微分形、物の中の形、エネルギーの流れ、動く人から見た場、数値で解く方法を、数学 8-9〜8-14 の道具で調べます。

\noindent この章の記号：面積を $A$ と書きます。⑦で、エネルギーの流れの矢印を $\vec S$ と書くためです（第1部・第2部では面積を $S$ と書きました）。閉じた面の積分のタイルも $\vec n\,\dd A$ と書きます。8-11・8-14 の $\vec n\,\dd S$ と同じものです。

%% ======================== 物理の道具 ========================
\bksection{doguu}{物理の道具}

\phstage{me}{法則を並べて、ほころびを探す}{3}

\bksubsection{1}{これまでの法則を4つ並べる}

\leavevmode\label{ph:電-14:sec1}第1部・第2部で学んだ電場と磁場の法則を、8-14①と同じ順に並べます。どれも「閉じた袋」か「閉じた道」を選んで、その上で場をたし合わせる形（\textbf{積分形}）です。

\begin{bkcenter}
\begin{bktable}{colspec={Q[c,m,wd=8mm] Q[c,m,wd=50mm] Q[l,m,wd=40mm] Q[c,m,wd=20mm]}, rowsep=3pt, colsep=2pt, cells={font=\small}, row{1}={font=\sffamily\bfseries\small, bg=bkTheme!10}}
 & 積分形（袋・道で） & 読み方 & 学んだ所 \\
（1） & $\displaystyle\oiint\vec E\cdot\vec n\,\dd A=\frac{Q_{\text{中}}}{\varepsilon_0}$ & 電荷は電場のわき出し（ガウスの法則） & 電-2 ③ \\
（2） & $\displaystyle\oiint\vec B\cdot\vec n\,\dd A=0$ & 磁場にわき出しはない & 電-6 ⑦ \\
（3） & $\displaystyle\oint\vec E\cdot\dd\vec r=-\frac{\dd\Phi_B}{\dd t}$ & 変わる磁場は電場の渦を作る（ファラデーの法則） & 電-8 ③ \\
（4） & $\displaystyle\oint\vec B\cdot\dd\vec r=\mu_0I_{\text{中}}$ & 電流は磁場の渦を作る（アンペールの法則） & 電-6 ⑥ \\
\end{bktable}
\bktabcaption{第1部・第2部で学んだ4つの法則（真空の中。電荷や電流はあってよい）。$Q_{\text{中}}$ は袋の中の電荷、$I_{\text{中}}$ は道を通りぬける電流、$\Phi_B$ は道をふちとする面を通る磁束。（4）は、まだ8-14①の表と同じ形ではない。}
\label{ph:電-14:tab-yottsu}
\end{bkcenter}

\begin{bkfigure}
\begin{tikzpicture}[x=1mm, y=1mm]
  \foreach \cx/\cy/\tt in {0/0/{（1）電荷は、電場のわき出し}, 66/0/{（2）磁力線は、いつも輪になる}, 0/-40/{（3）変わる磁場が、電場の渦を作る}, 66/-40/{（4）電流が、磁場の渦を作る}} {
    \draw[bkLine, line width=0.6pt, rounded corners=2mm, fill=white] (\cx,\cy) rectangle ++(62,-36);
    \node[anchor=north west, sxo lab, font=\sffamily\bfseries\footnotesize] at (\cx+2.5,\cy-2) {\tt};}
  % (1)
  \begin{scope}[shift={(31,-22)}]
    \filldraw[fill=bkRed!70, draw=bkRed!80!black] (0,0) circle (2.2);\node[text=white, font=\sffamily\bfseries\footnotesize] at (0,0) {＋};
    \foreach \a in {0,45,...,315} {\draw[sxo E] (\a:3.6) -- (\a:11);}
  \end{scope}
  % (2)
  \begin{scope}[shift={(97,-22)}]
    \foreach \s in {4,7,10} {\draw[bkBlue!80!black, line width=0.8pt, dash pattern=on 2.4pt off 1.4pt] (0,0) ellipse ({\s+4} and {0.8*\s});}
    \fill[bkRed!75!black] (-5,-1.8) rectangle (0,1.8);\fill[bkBlue!70!black] (0,-1.8) rectangle (5,1.8);
    \node[text=white, font=\sffamily\bfseries\footnotesize] at (-2.5,0) {N};\node[text=white, font=\sffamily\bfseries\footnotesize] at (2.5,0) {S};
  \end{scope}
  % (3)
  \begin{scope}[shift={(25,-64)}]
    \foreach \x in {-3,0,3} {\draw[sxo B, line width=1.1pt] (\x,-9) -- (\x,9);}
    \node[sxo lab, anchor=west, text=bkBlue!80!black] at (3.6,9.6) {$\vec B$ が強くなる};
    \draw[bkRed!80!black, line width=1.2pt] (0,0) ellipse (9 and 3);
    \draw[sxo E, line width=1.2pt] (-4.5,-2.598) -- (-6,-2.245);
    \node[sxo lab, anchor=east, text=bkRed!80!black] at (-9.6,-1) {$\vec E$ の渦};
  \end{scope}
  % (4)
  \begin{scope}[shift={(93,-64)}]
    \draw[bkOrange!80!black, line width=2pt] (0,-10) -- (0,10);
    \draw[sxo I, bkOrange!80!black, line width=1.4pt] (0,4) -- (0,10.5);
    \node[sxo lab, anchor=west] at (1,10) {電流 $I$};
    \foreach \s in {4,7.5} {\draw[bkBlue!80!black, line width=0.9pt, dash pattern=on 2.4pt off 1.4pt] (0,-2) ellipse [x radius=\s, y radius={0.35*\s}];}
    \draw[sxo B, line width=1.1pt] (3.2,-4.374) -- (4.6,-4.072);
    \node[sxo lab, anchor=west, text=bkBlue!80!black] at (8,-6) {$\vec B$ の渦};
  \end{scope}
\end{tikzpicture}
\bkfigcaption{4つの法則が言っていること。電場は実線の矢印、磁場は破線で区別する。（4）の磁場は、電流の向きに右ねじを進めるときにねじを回す向きに回る。}
\label{ph:電-14:fig-yottsu}
\end{bkfigure}

8-14①の表と比べると、（4）だけが違います。8-14①の（4）には、右辺にもう1つの項 $\mu_0\varepsilon_0\,\dd\Phi_E/\dd t$ がありました。この項は、どこから来たのでしょう。②で、その必要を実験の場面で確かめます。

\begin{phmijika}[4つの式は、身の回りのどこで働いている？]
（1）は、静電気で髪の毛が下じきに引かれるときの電場を決めます。（2）は、磁石をいくら細かく切っても、N極だけ・S極だけの磁石ができないことです。（3）は、発電機・変圧器・IH 調理器（電-8）。（4）は、電磁石・モーター（電-6・電-12）。私たちの家の電気製品は、ほとんどがこの4つの式の上で動いています。
\end{phmijika}

\bksubsection{2}{コンデンサーを充電すると、アンペールの法則が困る}

\leavevmode\label{ph:電-14:sec2}円い2枚の極板のコンデンサーを、導線で電池につないで充電します。充電している間、導線には電流 $I$ が流れ、極板に電荷がたまっていきます。導線の周りには、電流の磁場ができます。

では、導線を囲む円い道 $C$ に、アンペールの法則（4）を使ってみます。（4）の「道を通りぬける電流」は、道 $C$ をふちとする\textbf{面}を、電流がつらぬく量です。ふちが同じなら、面はどんな形に選んでもよいはずでした（8-11⑬）。ところが、図\ref{ph:電-14:fig-futatsu}の2つの面で、答えが違ってしまいます。

\begin{bkfigure}
\begin{tikzpicture}[x=1mm, y=1mm]
  % 導線
  \draw[bkOrange!80!black, line width=1.6pt] (0,0) -- (52,0) (68,0) -- (120,0);
  \sxocap{(60,0)}{14}{6}
  \node[sxo lab, anchor=east] at (51.6,7) {＋$Q$};\node[sxo lab, anchor=west] at (68.4,7) {－$Q$};
  \draw[sxo I, bkOrange!80!black] (6,2.2) -- (16,2.2); \node[sxo lab, anchor=south] at (11,3) {電流 $I$};
  \draw[sxo I, bkOrange!80!black] (104,2.2) -- (114,2.2); \node[sxo lab, anchor=south] at (109,3) {電流 $I$};
  % 極板の間の電場
  \foreach \yy in {-10,-5,0,5,10} {\draw[sxo E] (55,\yy) -- (65,\yy);}
  % 道 C
  \draw[bkInk, line width=1.0pt] (36,0) ellipse (2.5 and 9);
  \draw[-{Stealth[length=2.4pt]}, bkInk, line width=1.0pt] (33.6,-2.6) -- (33.55,-1.4);
  \node[sxo lab, anchor=north east] at (33,-8.5) {道 $C$};
  % 面1：平らな円板
  \fill[bkTheme!25, opacity=0.6] (36,0) ellipse (2.5 and 9);
  \draw[bkTheme!80!black, line width=0.5pt] (31,20) -- (35.5,7);
  \node[sxo lab, anchor=south] at (28,20) {面1（平らな円板）};
  % 面2：袋
  \draw[bkPurple!80!black, line width=0.9pt, dash pattern=on 3pt off 1.5pt] (36,9) .. controls (50,20) and (58,20) .. (60,18) .. controls (61,10) and (61,-10) .. (60,-18) .. controls (58,-20) and (50,-20) .. (36,-9);
  \node[sxo lab, anchor=south west, text=bkPurple!80!black] at (62,18) {面2（極板の間を通る袋）};
  \node[sxo lab, anchor=north, align=center] at (60,-21) {極板の間：電流はないが、\\電場 $\vec E$ が強くなっていく};
\end{tikzpicture}
\bkfigcaption{充電しているコンデンサーと、同じ道 $C$ をふちとする2つの面。面1は導線につらぬかれる。面2（破線）は極板の間を通るので、どの導線にもつらぬかれない。}
\label{ph:電-14:fig-futatsu}
\end{bkfigure}

\begin{itemize}[leftmargin=1.4\zw, itemsep=1pt, topsep=2pt, label={\color{bkTheme}●}]
\item 面1（道 $C$ に張った平らな円板）：導線がつらぬくので、通りぬける電流は $I$。（4）から、道 $C$ の上の磁場は $B=\mu_0I/(2\pi r)$。
\item 面2（同じ道 $C$ から、袋のようにふくらませて、極板の間を通す面）：極板の間には電荷が流れていないので、通りぬける電流は0。（4）から、道 $C$ の上の磁場は0。
\end{itemize}

同じ道 $C$ の上の同じ磁場が、面の選び方で $\mu_0I/(2\pi r)$ にも0にもなる。これは、どちらかがまちがっていることを示しています。実験すると、導線の周りにはちゃんと磁場があります。まちがっているのは、面2の計算、つまり（4）の式のほうです。

\begin{phmijika}[水道管で考える]
水道管の途中に、ゴムの膜で仕切った風船のような部屋があるとします。左から水を押しこむと、膜がふくらんで、右の水が押し出されます。膜を通りぬける水はないのに、管の右側でも水が流れます。コンデンサーも同じで、極板の間を電荷は通りぬけませんが、回路のどこでも同じ電流が流れます。膜の「ふくらみの変化」にあたるのが、極板の間の\textbf{電場の変化}です。
\end{phmijika}

\phstage{te}{変わる電場を、電流として数える}{3}

\bksubsection{3}{変位電流：変わる電場も、磁場を作る}

\leavevmode\label{ph:電-14:sec3}面2を通るものは、電流ではなく電場です。しかも、充電している間、その電場は強くなっていきます。そこで、極板の間の電場の流束 $\Phi_E$（面を通りぬける電場の量。電-2 ③）が、どれだけの速さで変わるかを数えてみます。

極板の面積を $A$、すき間を $d$、極板の電荷を $Q$ とします。極板の間の電場は一様で、$E=V/d$ でした（電-3 ④）。$Q=CV$、$C=\varepsilon_0A/d$ を使うと
\[ E=\frac{V}{d}=\frac{Q}{Cd}=\frac{Q}{\varepsilon_0A} \]
です。極板の間を通りぬける電場の流束は、電場×面積なので
\pheq{\Phi_E=EA=\frac{Q}{\varepsilon_0}}{V·m ＝ C ÷ (F/m)}
です。これは、ガウスの法則（1）で＋の極板を袋で囲んだときと同じ答えです。電荷 $Q$ は、1秒に $I$ ずつ増えていきます（$\dd Q/\dd t=I$。電-4 ①）。だから、$\Phi_E$ が増える速さは
\[ \frac{\dd\Phi_E}{\dd t}=\frac{1}{\varepsilon_0}\,\frac{\dd Q}{\dd t}=\frac{I}{\varepsilon_0} \]
です。両辺に $\varepsilon_0$ をかけると、ちょうど導線の電流 $I$ になります。そこで
\pheq{I_{\mathrm d}=\varepsilon_0\,\frac{\dd\Phi_E}{\dd t}}{A ＝ (F/m) × (V·m/s)}
\label{ph:電-14:hen}を\textbf{変位電流}とよびます。\bkidx{へんいでんりゅう}{変位電流}電荷は動いていませんが、電場の変化を「電流のなかま」として数えるのです。マクスウェルは、この変位電流も、ふつうの電流と同じように磁場の渦を作ると考え、アンペールの法則を次のように直しました。
\pheq{\oint\vec B\cdot\dd\vec r=\mu_0\Bigl(I_{\text{中}}+\varepsilon_0\frac{\dd\Phi_E}{\dd t}\Bigr)}{T·m ＝ (N/A²) × A}
\label{ph:電-14:amp-max}これを\textbf{アンペール・マクスウェルの法則}といいます。\bkidx{あんぺーる・まくすうぇるのほうそく}{アンペール・マクスウェルの法則}こうすると、面1では「電流 $I$＋変位電流0」、面2では「電流0＋変位電流 $I$」で、どちらの面でも同じ $\mu_0I$ になり、つじつまが合います。

\begin{phyomi}[$I_{\mathrm d}=\varepsilon_0\,\dd\Phi_E/\dd t$]
「電場の流束が1秒にどれだけ増えるか、に $\varepsilon_0$ をかけたものは、電流と同じはたらきをする」。電流と同じ単位 A です。電場が変わらなければ、変位電流は0です。だから、電池をつないだまま十分時間がたって充電が終わったコンデンサー（電-5 ⑪）や、止まった電荷の周りでは、アンペールの法則はもとの形のままです。
\end{phyomi}

\begin{bkfigure}
\begin{tikzpicture}
\begin{axis}[bkplot, width=100mm, height=50mm, xmin=0, xmax=10, ymin=0, ymax=9, xtick={0,2,4,5,6,8,10}, ytick={0,2,4,6,8}, xlabel={中心軸からの距離 $r$（cm）}, ylabel={$B$（$10^{-6}$\,T）}]
  \addplot[bkcurve, domain=0:5, samples=2] {0.8*x};
  \addplot[bkcurve, domain=5:10, samples=60] {20/x};
  \addplot[bkcurve2, dashed, domain=0:5, samples=2] {1.6*x};
  \addplot[bkcurve2, dashed, domain=5:10, samples=60] {40/x};
  \draw[bkGray, densely dotted] (axis cs:5,0) -- (axis cs:5,9);
  \node[font=\sffamily\footnotesize, anchor=west, align=left, inner sep=1.5pt] at (axis cs:6.2,7.6) {破線：電流2倍（2 A）};
  \node[font=\sffamily\footnotesize, anchor=west, align=left, inner sep=1.5pt] at (axis cs:7.0,1.0) {実線：電流1 A};
  \node[font=\sffamily\footnotesize, anchor=south east, inner sep=1pt] at (axis cs:4.8,0.3) {極板の縁 $R=\qty{5}{cm}$};
\end{axis}
\end{tikzpicture}
\bkfigcaption{半径 \qty{5}{cm} の円い極板を充電するときの、極板の間の磁場の大きさ（例題1・11）。極板の内側では、道が囲む変位電流が $r^2$ に比例するので、$B$ は $r$ に比例する。外側では導線の周りと同じ $\mu_0I/(2\pi r)$。}
\label{ph:電-14:fig-Br}
\end{bkfigure}

\begin{phbaisuru}{充電の電流 $I$}
変位電流も2倍になり、極板の間の磁場も、どこでも2倍になります（図\ref{ph:電-14:fig-Br}の破線）。極板の内側の、中心から $r$ の円い道が囲む変位電流は、全体の $(r/R)^2$ です（電場は一様で、流束は面積に比例するから）。$2\pi r\times B=\mu_0I(r/R)^2$ から $B=\mu_0Ir/(2\pi R^2)$ で、中心では0、縁でいちばん強くなります（例題1）。
\end{phbaisuru}

\Needspace{12\baselineskip}
\begin{motto}[4]
変位電流の名前は、マクスウェルが、真空の中にも何か「ずれる（変位する）もの」があると考えたことから来ています。今では、真空の中には動く物はなく、変わる電場そのものが磁場を作る、と考えます。誘電体（ガラスやプラスチック）の中では、分子の中の電荷が本当に少しずれる（分極：電-3 ⑥）ので、その分も変位電流にふくまれます（⑥）。
\end{motto}

\bksubsection{4}{電荷の保存と、つじつまを合わせる}

\leavevmode\label{ph:電-14:sec4}変位電流は、つじつまを合わせるために付け加えた「おまけ」なのでしょうか。そうではありません。電荷の保存（電荷は勝手に生まれも消えもしない：電-4 ⑥）と、深く結びついています。

\begin{bkfigure}
\begin{tikzpicture}[x=1mm, y=1mm]
  \draw[bkOrange!80!black, line width=1.6pt] (0,0) -- (40,0) (60,0) -- (100,0);
  \sxocap{(50,0)}{14}{6}
  \draw[sxo I, bkOrange!80!black] (14,2.2) -- (26,2.2); \node[sxo lab, anchor=south] at (20,3) {入る電流 $I$};
  \foreach \yy in {-10,-5,0,5,10} {\draw[sxo E] (45,\yy) -- (55,\yy);}
  \draw[bkPurple!80!black, line width=1.0pt, dash pattern=on 3pt off 1.5pt, rounded corners=4mm] (30,-19) rectangle (50.8,19);
  \node[sxo lab, anchor=south, text=bkPurple!80!black] at (40,19.6) {閉じた袋};
  \node[sxo lab, anchor=west, align=left] at (62,13) {袋の右の面を\\電場が出ていく};
  \draw[bkGray, line width=0.4pt] (61.6,11.5) -- (53,8);
  \node[sxo lab, anchor=north, align=center] at (40,-20) {＋の極板を囲む};
\end{tikzpicture}
\bkfigcaption{＋の極板を囲む閉じた袋。左の面から電流 $I$ が入り、右の面（極板の間）から電場が出ていく。}
\label{ph:電-14:fig-fukuro}
\end{bkfigure}

＋の極板を、閉じた袋で囲みます（図\ref{ph:電-14:fig-fukuro}）。袋には、導線から電流 $I$ が入ってきて、出ていく電流はありません。だから、袋の中の電荷は $\dd Q_{\text{中}}/\dd t=I$ の速さで増えます。一方、ガウスの法則（1）から、袋から出ていく電場の流束は $\Phi_E=Q_{\text{中}}/\varepsilon_0$ です。時間で微分して $\varepsilon_0$ をかけると
\[ \varepsilon_0\,\frac{\dd\Phi_E}{\dd t}=\frac{\dd Q_{\text{中}}}{\dd t}=I \]
です。「袋に入る電流」と「袋から出る変位電流」が、ちょうど等しいのです。\textbf{電流と変位電流をたしたもの}は、どんな袋に対しても、入る量と出る量が等しく、どこかで生まれたり消えたりしません。だから、どんな面を選んでも、アンペール・マクスウェルの法則の右辺は同じ値になります。変位電流は、電荷の保存とガウスの法則から、どうしても必要になる項なのです。微分形で同じことを確かめたのが 8-14④ です。

\begin{phmijika}[交流がコンデンサーを「通る」わけ]
交流の回路では、コンデンサーをつないでも電流が流れ続けます（電-10 ④）。電荷は極板の間を通りぬけていませんが、極板の間では電場が行ったり来たりして、変位電流が流れています。電流と変位電流をつなげると、回路を1周する「電流」は、とぎれずにつながっています。
\end{phmijika}

\phkokomade

\Needspace{14\baselineskip}
\phstage{shiki}{場所ごとの式とエネルギーで読む}{5}

\bksubsection{5}{積分形と微分形：マクスウェル方程式}

\leavevmode\label{ph:電-14:sec5}アンペール・マクスウェルの法則を（4）に入れると、4つの式がそろいます。これが\textbf{マクスウェル方程式}です。\bkidx{まくすうぇるほうていしき}{マクスウェル方程式}

\begin{bkcenter}
\begin{bktable}{colspec={Q[c,m,wd=8mm] Q[c,m,wd=58mm] Q[c,m,wd=48mm]}, rowsep=3pt, colsep=2pt, cells={font=\small}, row{1}={font=\sffamily\bfseries\small, bg=bkTheme!10}}
 & 積分形（袋・道で） & 微分形（場所ごとに） \\
（1） & $\displaystyle\oiint\vec E\cdot\vec n\,\dd A=\frac{Q_{\text{中}}}{\varepsilon_0}$ & $\mathrm{div}\,\vec E=\dfrac{\rho}{\varepsilon_0}$ \\
（2） & $\displaystyle\oiint\vec B\cdot\vec n\,\dd A=0$ & $\mathrm{div}\,\vec B=0$ \\
（3） & $\displaystyle\oint\vec E\cdot\dd\vec r=-\frac{\dd\Phi_B}{\dd t}$ & $\mathrm{rot}\,\vec E=-\dfrac{\partial\vec B}{\partial t}$ \\
（4） & $\displaystyle\oint\vec B\cdot\dd\vec r=\mu_0I_{\text{中}}+\mu_0\varepsilon_0\frac{\dd\Phi_E}{\dd t}$ & $\mathrm{rot}\,\vec B=\mu_0\vec j+\mu_0\varepsilon_0\dfrac{\partial\vec E}{\partial t}$ \\
\end{bktable}
\bktabcaption{マクスウェル方程式（真空の中。電荷や電流はあってよい）。8-14①の表と同じ。$\rho$ は電荷密度（\qty{1}{m^3} あたりの電荷）、$\vec j$ は電流密度（\qty{1}{m^2} あたりの電流）。}
\label{ph:電-14:tab-maxwell}
\end{bkcenter}

左の積分形と右の微分形は、ガウスの定理・ストークスの定理でつながっていて、同じことを言っています（8-14③で、（3）を例に確かめた）。積分形は、対称な形の場の大きさを求めるのが得意です（電-2 ④の球・直線・平面、電-6 ⑥のソレノイド）。微分形は、場所ごとに何が起きているかを見るのが得意で、電磁波の式（電-15 ⑥）や、コンピュータで場を計算する方法（⑩）は、微分形から作ります。

\Needspace{10\baselineskip}
\begin{phflow}[時間で変わらないとき]
\phflowrow{（3）$\displaystyle\oint\vec E\cdot\dd\vec r=-\frac{\dd\Phi_B}{\dd t}$}{磁場が\\変わらない}{$\displaystyle\oint\vec E\cdot\dd\vec r=0$\\電位が使える（電-2 ⑨）}
\phflowrow{（4）$\displaystyle\oint\vec B\cdot\dd\vec r=\mu_0I_{\text{中}}+\mu_0\varepsilon_0\frac{\dd\Phi_E}{\dd t}$}{電場が\\変わらない}{$\displaystyle\oint\vec B\cdot\dd\vec r=\mu_0I_{\text{中}}$\\アンペールの法則（電-6 ⑥）}
\end{phflow}

場が時間で変わらないとき、（1）（3）は電場だけ、（2）（4）は磁場だけの式に分かれます。電場と磁場は、別々に考えてよいのです。第1部（静電場）と第2部のはじめ（電流の磁場）を、別々に学べたのはこのためです。電場と磁場が結びつくのは、（3）の $\dd\Phi_B/\dd t$ と（4）の $\dd\Phi_E/\dd t$、つまり\textbf{場が時間で変わるとき}だけです。

\begin{phmada}[N極だけの磁石（磁気単極子）はあるの？]
（2）の $\mathrm{div}\,\vec B=0$ は、「N極だけ・S極だけの粒（\textbf{磁気単極子}）は、ない」と言っています。\bkidx{じきたんきょくし}{磁気単極子}ところが、1931年にディラックが、もし磁気単極子が1つでも宇宙にあれば、電気素量 $e$ がなぜ決まった値なのかを説明できることを示しました。その後、素粒子の理論の多くも、宇宙のはじめに磁気単極子ができた可能性を予想しています。

そこで、宇宙線の観測や、加速器の実験で、磁気単極子が探されてきました。2018年には、スイスの大型加速器（LHC）で鉛の原子核どうしをぶつけて、宇宙でいちばん強いと考えられる磁場を作り、その中で磁気単極子が生まれないかを調べる実験（MoEDAL 実験）が行われましたが、見つかりませんでした。今のところ、磁気単極子は1つも確かに見つかっていません。見つかれば、（2）の右辺は0でなくなり、マクスウェル方程式の形が変わります。
\end{phmada}

\bksubsection{6}{物の中の形：$\vec D$ と $\vec H$ を使う4つの式}

\leavevmode\label{ph:電-14:sec6}ガラスや水のような誘電体の中では、電場にひっぱられて分子の中の電荷が少しずれ（分極）、表面にも中にも、ずれてできた電荷が現れます（電-3 ⑥）。鉄のような磁性体の中では、原子の小さな磁石がそろって（磁化）、磁化の電流のようなものが現れます（電-7 ③）。こうした「物の中で生まれる電荷・電流」を、いちいち数えるのはたいへんです。

そこで、電-3 ⑧ では分極の分を入れた\textbf{電束密度} $\vec D=\varepsilon_0\vec E+\vec P$ を、電-7 ③ では磁化の分を差し引いた\textbf{磁場の強さ} $\vec H=\vec B/\mu_0-\vec M$ を作りました。これを使うと、自分で決めて置いた電荷（\textbf{自由電荷}）と、導線を流れる電流（\textbf{自由電流}）だけで書けます。

\begin{bkcenter}
\begin{bktable}{colspec={Q[c,m,wd=8mm] Q[c,m,wd=62mm] Q[c,m,wd=46mm]}, rowsep=3pt, colsep=2pt, cells={font=\small}, row{1}={font=\sffamily\bfseries\small, bg=bkTheme!10}}
 & 積分形 & 微分形 \\
（1） & $\displaystyle\oiint\vec D\cdot\vec n\,\dd A=Q_{\text{自由}}$ & $\mathrm{div}\,\vec D=\rho_{\text{自由}}$ \\
（2） & $\displaystyle\oiint\vec B\cdot\vec n\,\dd A=0$ & $\mathrm{div}\,\vec B=0$ \\
（3） & $\displaystyle\oint\vec E\cdot\dd\vec r=-\frac{\dd\Phi_B}{\dd t}$ & $\mathrm{rot}\,\vec E=-\dfrac{\partial\vec B}{\partial t}$ \\
（4） & $\displaystyle\oint\vec H\cdot\dd\vec r=I_{\text{自由}}+\frac{\dd\Phi_D}{\dd t}$ &$\mathrm{rot}\,\vec H=\vec j_{\text{自由}}+\dfrac{\partial\vec D}{\partial t}$ \\
\end{bktable}
\bktabcaption{物の中のマクスウェル方程式。（4）の $\Phi_D$ は、道をふちとする面を通る $\vec D$ の流束で、その変わる速さが物の中の変位電流。（2）（3）は物の中でも同じ形。}
\label{ph:電-14:tab-mono}
\end{bkcenter}

ふつうの物（電場や磁場がそれほど強くなく、時間で性質が変わらない物）では、$\vec D=\varepsilon\vec E$（$\varepsilon=\varepsilon_{\mathrm r}\varepsilon_0$）、$\vec B=\mu\vec H$（$\mu=\mu_{\mathrm r}\mu_0$）です。真空の式の $\varepsilon_0$ を $\varepsilon$ に、$\mu_0$ を $\mu$ にかえればよいのです。たとえば水（$\varepsilon_{\mathrm r}\approx80$）の中では、同じ電荷の周りの電場が、真空の80分の1になります（例題10）。

\begin{phmodel}[物の中のマクスウェル方程式]
\phignore{物の分子1つ1つの周りの細かな場（原子の大きさの、とがった場）。これを、たくさんの分子でならした平均の場を $\vec E$・$\vec B$ と考える。}
\phassume{$\vec D=\varepsilon\vec E$、$\vec B=\mu\vec H$ で、$\varepsilon$・$\mu$ は場の強さによらない一定の数（線形）。物の性質は時間で変わらない。}
\phlimit{とても強い電場（絶縁が破れる）。鉄のように、$\vec B$ と $\vec H$ が比例せず、前の磁化をおぼえている物（ヒステリシス：電-7 ⑧）。場が速く変わり、$\varepsilon$ が周波数で変わるとき（電-15 ⑨の分散）。}
\end{phmodel}

\bksubsection{7}{電磁場のエネルギーと、その流れ：ポインティングベクトル}

\leavevmode\label{ph:電-14:sec7}電場と磁場には、エネルギーがたくわえられています。コンデンサーの極板の間の電場には \qty{1}{m^3} あたり $\tfrac12\varepsilon_0E^2$（電-3 ⑦）、コイルの中の磁場には \qty{1}{m^3} あたり $B^2/(2\mu_0)$ のエネルギーがありました（電-8 ⑧）。2つを合わせた
\pheq{u=\frac12\varepsilon_0E^2+\frac{B^2}{2\mu_0}}{J/m³}
\label{ph:電-14:u}を、電磁場の\textbf{エネルギー密度}といいます。\bkidx{えねるぎーみつど}{エネルギー密度}

場が時間で変わると、ある場所のエネルギーが増えたり減ったりします。エネルギーは保存するので、減った分は、どこかへ\textbf{流れて}いったか、物を温める仕事（ジュール熱）に使われたはずです。1884年、イギリスのポインティングは、マクスウェル方程式から、そのエネルギーの流れを表す矢印を見つけました。
\pheq{\vec S=\vec E\times\vec H}{W/m² ＝ (V/m) × (A/m)}
\label{ph:電-14:poynting}この矢印 $\vec S$ を\textbf{ポインティングベクトル}といいます。\bkidx{ぽいんてぃんぐべくとる}{ポインティングベクトル}真空の中では $\vec H=\vec B/\mu_0$ です。

\begin{phyomi}[$\vec S=\vec E\times\vec H$]
「電場と磁場が両方ある所では、エネルギーが、2つに直角な向きに流れる。その強さは、\qty{1}{m^2} の窓を1秒に通るエネルギー（W/m²）」。向きは外積（4-13）の向きで、右手の親指以外の指を $\vec E$ から $\vec H$ へ回すと、親指が $\vec S$ の向きです。電場だけ、または磁場だけの所では、エネルギーは流れません。止まった電荷と磁石を並べると、両方の場がある所で $\vec S$ は0ではありませんが、輪になってぐるぐる回っているだけで、どこにもたまらず、どこからも出ていきません。
\end{phyomi}

ポインティングが示したエネルギーの保存の式は、閉じた袋について、次の形です。
\pheq{-\frac{\dd U_{\text{中}}}{\dd t}=\oiint\vec S\cdot\vec n\,\dd A+P_{\text{熱}}}{W}
\label{ph:電-14:hozon}$U_{\text{中}}$ は袋の中の電磁場のエネルギー（$u$ を袋の中でたしたもの）、$P_{\text{熱}}$ は袋の中で電流が物にする仕事率（ジュール熱など）です。「袋の中の電磁場のエネルギーが減る速さ＝袋から外へ流れ出る速さ＋熱になる速さ」と読みます。熱-2 のエネルギーの保存（第一法則）や、力-13 の流れのエネルギーの式と同じ形です。この巻の前付の「保存則と相似の地図」で、エネルギーの行の電磁気の欄に書いた「ポインティングベクトル」は、この式のことです。

\begin{phmijika}[コンデンサーに入るエネルギーは、横から入る]
充電しているコンデンサーの極板の間では、電場 $\vec E$（＋から－へ）と、変位電流が作る磁場（極板の縁にそって回る）が両方あります。外積の向きを調べると、$\vec S$ は、極板の縁から\textbf{中心へ向かって}います。コンデンサーにたまるエネルギーは、導線を通って極板に入るのではなく、すき間の横の空間から流れこむのです（例題7）。
\end{phmijika}

\begin{bkfigure}
\begin{tikzpicture}[x=1mm, y=1mm]
  % 側面図
  \fill[bkGray!70] (14,30) rectangle (56,32);\fill[bkGray!70] (14,0) rectangle (56,2);
  \node[sxo lab, anchor=south] at (35,32.6) {＋の極板};\node[sxo lab, anchor=north] at (35,-0.6) {－の極板};
  \foreach \x in {20,30,40,50} {\draw[sxo E] (\x,28) -- (\x,4);}
  \node[sxo lab, text=bkRed!80!black, anchor=west] at (50.6,24) {$\vec E$};
  \draw[sxo S] (71,22) -- (59,22); \draw[sxo S] (71,10) -- (59,10);
  \draw[sxo S] (-1,22) -- (11,22); \draw[sxo S] (-1,10) -- (11,10);
  \node[sxo lab, text=bkOrange!75!black, anchor=south] at (65,23.5) {$\vec S$};
  \node[sxo lab, anchor=north] at (35,-5) {横から見た図};
  % 上から見た図
  \begin{scope}[shift={(102,16)}]
    \draw[bkGray!80!black, line width=0.8pt, fill=bkGray!15] (0,0) circle (11);
    \foreach \a in {0,45,...,315} {\draw[sxo S, line width=1.1pt] (\a:20) -- (\a:12);}
    \draw[bkBlue!80!black, line width=0.9pt, dash pattern=on 2.4pt off 1.4pt] (0,0) circle (15);
    \draw[sxo B] (15,0.2) -- (15,2);
    \node[sxo lab, text=bkBlue!80!black, anchor=west] at (16.8,-6.6) {$\vec H$};
    \node[sxo lab, align=center] at (0,0) {$\vec E$ は\\奥へ};
    \node[sxo lab, anchor=north] at (0,-21) {上から見た図};
  \end{scope}
\end{tikzpicture}
\bkfigcaption{充電しているコンデンサーに流れこむエネルギー。極板の間の電場（実線の矢印）と、その周りを回る磁場（破線）から、ポインティングベクトル（太い矢印）は、縁から中心へ向かう。}
\label{ph:電-14:fig-condS}
\end{bkfigure}

\begin{phbaisuru}{電場 $E$}
電場のエネルギー密度 $\tfrac12\varepsilon_0E^2$ は、$E$ の2乗に比例するので4倍になります。同じ磁場のもとでは、ポインティングベクトルは $E$ に比例するので2倍です。エネルギーの「たまり方」は2乗、「流れ方」は電場と磁場の積で決まります。
\end{phbaisuru}

\bksubsection{8}{電線の周りを流れるエネルギー}

\leavevmode\label{ph:電-14:sec8}やってみようの答えです。電池から豆電球へ、エネルギーはどこを通るのでしょう。

\Needspace{5\baselineskip}\noindent\textbf{抵抗線に流れこむエネルギー}\par\nobreak
長さ $l$・半径 $a$ の抵抗線（ドライヤーのヒーター線など）の両はしに電圧 $V$ をかけ、電流 $I$ が流れているとします。抵抗線の表面では、電場は線にそった向きで、大きさは $E=V/l$ です（電-4 ③）。磁場の強さは、アンペールの法則から、線を囲む円の上で $H=I/(2\pi a)$ です（電-6 ⑤。$H=B/\mu_0$）。電場は線にそって、磁場は線を囲む向きなので、外積 $\vec S=\vec E\times\vec H$ は、表面に垂直に、\textbf{線の中へ向かって}います（図\ref{ph:電-14:fig-teikou}）。

\begin{bkfigure}
\begin{tikzpicture}[x=1mm, y=1mm]
  \filldraw[fill=bkOrange!25, draw=bkInk, line width=0.7pt] (10,-5) rectangle (110,5);
  \draw[bkInk, line width=0.7pt] (10,0) ellipse (1.6 and 5);
  \draw[bkInk, line width=0.7pt] (110,-5) arc (-90:90:1.6 and 5);
  \node[sxo lab, fill=bkOrange!25] at (61,0) {抵抗線（半径 $a$、長さ $l$）};
  \foreach \x in {14,66} {\draw[sxo E] (\x,7.5) -- (\x+10,7.5);}
  \node[sxo lab, text=bkRed!80!black, anchor=south west] at (6,10.6) {$\vec E$（線にそって）};
  \foreach \x in {32,90} {\draw[bkBlue!80!black, line width=0.8pt, dash pattern=on 2.4pt off 1.4pt] (\x,0) ellipse (3 and 9);
    \draw[sxo B] (\x-3,1) -- (\x-3,-1);}
  \node[sxo lab, text=bkBlue!80!black, anchor=north] at (90,-10.5) {$\vec H$（線を囲む）};
  \foreach \x in {50,58} {\draw[sxo S] (\x,22) -- (\x,7); \draw[sxo S] (\x,-22) -- (\x,-7);}
  \node[sxo lab, text=bkOrange!75!black, anchor=west] at (61,19) {$\vec S$（線の中へ）};
  \draw[sxo I] (-4,0) -- (7.5,0); \node[sxo lab, anchor=south] at (1,1) {電流 $I$};
  \node[sxo lab, anchor=north east] at (9,-6.5) {電位が高い};\node[sxo lab, anchor=north west] at (111,-6.5) {低い};
\end{tikzpicture}
\bkfigcaption{抵抗線の周りのエネルギーの流れ。電場（実線）は線にそい、磁場（破線）は線を囲む。ポインティングベクトル（太い矢印）は、表面から線の中へ入り、そこで熱になる。}
\label{ph:電-14:fig-teikou}
\end{bkfigure}

表面での大きさは $S=EH$ です。表面の面積 $2\pi al$ をかけて、線に入るエネルギーを全部たすと
\[ S\times2\pi al=\frac{V}{l}\times\frac{I}{2\pi a}\times2\pi al=VI \]
となります。\label{ph:電-14:VI}電-4 ⑤ の電力 $P=VI$ が、場のことばで出てきました。抵抗線がジュール熱で出しているエネルギーは、ちょうど、周りの空間から表面を通って流れこんでいるのです。

\Needspace{10\baselineskip}
\begin{phflow}[場から、回路の電力へ]
\phflowrow{$\displaystyle\oiint\vec S\cdot\vec n\,\dd A$\\（袋に入るエネルギー）}{一様な抵抗線\\の表面}{$P=VI$\\（電-4 ⑤）}
\end{phflow}

\Needspace{5\baselineskip}\noindent\textbf{電池から豆電球へ}\par\nobreak
乾電池と豆電球の回路で、行きの導線と帰りの導線の間には、電位の差があるので、＋側の導線から－側の導線へ向かう電場があります。導線の周りには、電流の磁場があります。この2つの外積は、2本の導線の間の空間で、電池から豆電球の向きをさしています（図\ref{ph:電-14:fig-denchi}）。電池の化学のエネルギーは、電池の中で電磁場のエネルギーに変わり、2本の導線にはさまれた空間を流れて、豆電球のフィラメントに表面から入り、熱と光になります。導線は、エネルギーの「通り道」の形を決める「レール」の役目をしているのです。

\begin{bkfigure}
\begin{tikzpicture}[x=1mm, y=1mm]
  \draw[bkOrange!80!black, line width=1.2pt] (18,26) -- (100,26) (18,4) -- (100,4);
  \draw[bkOrange!80!black, line width=1.2pt] (18,26) -- (18,19) (18,11) -- (18,4) (100,26) -- (100,19) (100,11) -- (100,4);
  \begin{scope}[shift={(18,15)}, rotate=-90]\sxobat{(0,0)}\end{scope}
  \filldraw[fill=bkYellow!50, draw=bkInk, line width=0.7pt] (100,15) circle (4);
  \draw[bkInk, line width=0.5pt] (97.5,13) .. controls (99,17) and (101,17) .. (102.5,13);
  \node[sxo lab, anchor=west] at (105,15) {豆電球};
  \node[sxo lab, anchor=east] at (12,15) {乾電池};
  \foreach \x in {34,58,82} {\draw[sxo E] (\x,24.5) -- (\x,5.5);}
  \node[sxo lab, text=bkRed!80!black, anchor=west] at (34.6,20) {$\vec E$};
  \foreach \x in {46,70} {\draw[sxo S] (\x-6,15) -- (\x+6,15);}
  \node[sxo lab, text=bkOrange!75!black, anchor=south] at (46,16.6) {$\vec S$};
  \node[sxo lab, anchor=south] at (60,27) {＋側の導線（$\vec H$ は導線を囲む）};
  \node[sxo lab, anchor=north] at (60,3) {－側の導線};
\end{tikzpicture}
\bkfigcaption{乾電池と豆電球の回路で、2本の導線の間を流れるエネルギー（真ん中の面の上だけをかいた）。電場（実線）は＋側から－側へ。ポインティングベクトル（太い矢印）は、電池から豆電球へ向かう。}
\label{ph:電-14:fig-denchi}
\end{bkfigure}

\Needspace{5\baselineskip}\noindent\textbf{同軸ケーブル}\par\nobreak
テレビのアンテナ線などの\textbf{同軸ケーブル}は、内側の導線（半径 $a$）を、外側の筒の形の導体（内側の半径 $b$）が囲み、その間をポリエチレンのような誘電体でうめたものです。内側の導線から外の導体へ向かう電場（中心から外向き）と、内側の導線の電流の磁場（中心を囲む向き）の外積は、ケーブルにそって、送る先へ向かいます。エネルギーは、2つの導体の\textbf{間の誘電体の中}を流れていきます（例題6）。外の導体の外側では、行きと帰りの電流の磁場が打ち消し合い、電場もないので、エネルギーは外へもれません。同軸ケーブルが、電波をもらさず、外の電波にもじゃまされにくいのはこのためです（電-13 ②）。

\begin{phmijika}[送電線のエネルギーは、空中を飛んでくる]
発電所から家まで、送電線の電線と電線の間、電線と地面の間の\textbf{空気の中}を、電磁場のエネルギーが流れています。電線そのものは、少しずつ熱になる分（電線の抵抗による損失：電-11 ⑥）を、周りから表面を通して受け取っているだけです。電気のエネルギーは「電線の中を流れる」と言うことが多いですが、場で見ると、電線は道案内なのです。
\end{phmijika}

\bksubsection{9}{電場と磁場は、見る人で入れかわる}

\leavevmode\label{ph:電-14:sec9}電-8 ④ では、一様な磁場 $\vec B$ の中を速さ $v$ で動く金属の棒に、起電力 $vBl$ が生まれることを、ローレンツ力（電-6 ③）で説明しました。この場面を、2人の人の目で見直します（図\ref{ph:電-14:fig-futari}）。

\begin{bkfigure}
\begin{tikzpicture}[x=1mm, y=1mm]
  \foreach \cx/\tt in {0/{止まっている人A}, 66/{棒といっしょに動く人B}} {
    \draw[bkLine, line width=0.6pt, rounded corners=2mm, fill=white] (\cx,0) rectangle ++(62,40);
    \node[anchor=north west, sxo lab, font=\sffamily\bfseries\footnotesize] at (\cx+2.5,38) {\tt};
    \foreach \x in {6,14,22} {\foreach \y in {6,14,22} {\node[text=bkBlue!80!black, font=\footnotesize] at (\cx+\x,\y) {$\times$};}}
    \filldraw[fill=bkGray!45, draw=bkInk, line width=0.6pt] (\cx+29,4) rectangle (\cx+33,27);
    \filldraw[fill=bkBlue!60, draw=bkInk, line width=0.4pt] (\cx+31,17) circle (1.4);
    \node[sxo lab, anchor=west] at (\cx+34,13) {電子};
  }
  \node[sxo lab, text=bkBlue!80!black, anchor=west] at (3,31) {$\vec B$（奥向き）};
  \draw[sxo I] (34,24) -- (44,24); \node[sxo lab, anchor=west] at (44.5,24) {速さ $v$};
  \draw[bkPurple!80!black, line width=1.1pt, -{Stealth[length=3pt]}] (31,15.2) -- (31,8.5);
  \node[sxo lab, anchor=west, text=bkPurple!80!black, align=left] at (34,5.5) {磁場の力 $q\vec v\times\vec B$};
  \draw[sxo E] (101,15) -- (101,26); \node[sxo lab, text=bkRed!80!black, anchor=west] at (102,22) {電場 $\vec E'$};
  \draw[bkPurple!80!black, line width=1.1pt, -{Stealth[length=3pt]}] (97,15.2) -- (97,8.5);
  \node[sxo lab, anchor=west, text=bkPurple!80!black, align=left] at (100,5.5) {電場の力 $q\vec E'$};
  \node[sxo lab, anchor=east, align=right] at (126,31) {棒も電子も止まって見える};
\end{tikzpicture}
\bkfigcaption{磁場の中を動く棒の中の電子。止まっている人 A には「磁場の力」、棒といっしょに動く人 B には「電場の力」に見える。どちらでも、電子は同じ向き・同じ大きさの力を受ける（電子は負の電荷なので、力は電場 $\vec E'$ と逆向き）。}
\label{ph:電-14:fig-futari}
\end{bkfigure}

\begin{itemize}[leftmargin=1.4\zw, itemsep=1pt, topsep=2pt, label={\color{bkTheme}●}]
\item 止まっている人 A：電子は棒といっしょに速さ $v$ で動いているので、磁場からローレンツ力 $q\vec v\times\vec B$ を受け、棒にそって動かされる。A の目には、電場はない。
\item 棒といっしょに動く人 B：棒も電子も止まって見える。止まった電荷には、磁場は力をはたらかせない。それでも電子は棒にそって動かされるので、B は「棒のある所に、電場 $\vec E'=\vec v\times\vec B$ がある」と考えるしかない。
\end{itemize}

A が「磁場」と見ているものの一部を、B は「電場」と見ているのです。電場と磁場は、それぞれが独立にあるものではなく、\textbf{1つの電磁場}を、見る人の動き方によって違う分け方で見たものです。
\pheq{\vec E'=\vec E+\vec v\times\vec B\qquad\text{（$v\ll c$ のとき）}}{V/m ＝ (m/s) × T}
\label{ph:電-14:Edash}\bkidx{でんじば}{電磁場}電-8 例題の棒（$v=\qty{2}{m/s}$、$B=\qty{0.1}{T}$）では、B の見る電場は \qty{0.2}{V/m} で、長さ \qty{0.5}{m} の棒の両はしの電圧は \qty{0.1}{V} です（例題8）。A がファラデーの法則やローレンツ力で計算した起電力と同じです。

逆に、止まった電荷のそばを動く人には、電荷が動いて見える（電流に見える）ので、磁場も見えます。速さが光の速さに近づくと、この入れかわりの式は、時間と空間の見え方が変わる\textbf{相対性理論}の式になります。1905年、アインシュタインの相対性理論の論文の題は「動いている物体の電気力学について」でした。電磁場の入れかわりを調べたことが、相対性理論の入口になったのです。そのくわしい話は、物理シリーズの『原子』（原-11）で学びます。

\Needspace{5\baselineskip}\noindent\textbf{はてな？の答え}\par\nobreak
\label{ph:電-14:hatena}\bkhatenaref に答えます。電気と磁気は、別々のものではありません。止まった電荷は電場だけを作りますが、動いている人から見ると、その電荷は電流になり、磁場も作ります。磁石のそばを動く人には、止まった人が「磁場」と見るものの一部が、「電場」に見えます。電場と磁場は、1つの「電磁場」の2つの顔なのです。そして、マクスウェル方程式の（3）と（4）が言うように、変わる磁場は電場を、変わる電場は磁場を作ります。この2つがかわるがわる生まれながら空間を進むものが、光や電波（電磁波）です（電-15）。磁石が紙切れを引かないのは、紙に電荷がほとんどなく、磁石のそばで止まっているからで、電気と磁気が別のものだからではありません。

\bksubsection{10}{数値で解く：FDTD 法の入口}

\leavevmode\label{ph:電-14:sec10}スマホのアンテナや、電子レンジの中の電磁場は、形が複雑で、式で解くことはできません。そこで、空間と時間を細かなマスに分けて、マクスウェル方程式を「となりのマスとの差」で計算します。いちばん広く使われているのが \textbf{FDTD 法}（時間領域の差分法）で、\bkidx{えふでぃーてぃーでぃーほう}{FDTD 法}1966年にイー（K. S. Yee）が考えました。8-13⑧ で弦の波を表計算で動かしたのと、同じ考えです。

\Needspace{5\baselineskip}\noindent\textbf{1次元の式}\par\nobreak
$x$ の向きに進む波で、電場が $y$ の向き（$E_y$）、磁場が $z$ の向き（$H_z$）だけのときを考えます。（3）を $x$-$y$ 面の小さな長方形の道に、（4）を $x$-$z$ 面の小さな長方形の道に使うと（8-14③ と同じ手順）、
\[ \frac{\partial H_z}{\partial t}=-\frac{1}{\mu_0}\,\frac{\partial E_y}{\partial x},\qquad \frac{\partial E_y}{\partial t}=-\frac{1}{\varepsilon_0}\,\frac{\partial H_z}{\partial x} \]
が出ます。「電場の場所による変化が磁場を変え、磁場の場所による変化が電場を変える」という、⑤の（3）（4）そのものです。2つの式から一方を消すと、8-13⑧ の波動方程式になり、速さ $c=1/\sqrt{\mu_0\varepsilon_0}$ で進む波が出てきます（電-15 ⑥）。

単位をそろえるため、磁場の強さに $\eta_0=\sqrt{\mu_0/\varepsilon_0}\approx\qty{377}{\ohm}$（電-15 ⑦）をかけた $h=\eta_0H_z$ を使うと、2つの式は $\partial h/\partial t=-c\,\partial E_y/\partial x$、$\partial E_y/\partial t=-c\,\partial h/\partial x$ と、同じ形になります。

\Needspace{5\baselineskip}\noindent\textbf{電場と磁場を、半マスずらして置く}\par\nobreak
イーの工夫は、電場を置く点と磁場を置く点を、空間でも時間でも\textbf{半マスずつずらす}ことでした（図\ref{ph:電-14:fig-yee}）。電場のとなりの2点の差で、その間の磁場を進め、磁場のとなりの2点の差で、その間の電場を進めます。$c\,\Delta t=\Delta x$ にとると、式はとても簡単になります。
\[ h_{i+\frac12}\leftarrow h_{i+\frac12}-\bigl(E_{i+1}-E_i\bigr),\qquad E_i\leftarrow E_i-\bigl(h_{i+\frac12}-h_{i-\frac12}\bigr) \]
（$\leftarrow$ は「右の値で、左を書きかえる」。8-13 の表計算の1歩と同じです。）

\begin{bkfigure}
\begin{tikzpicture}[x=1mm, y=1mm]
  \draw[bkGray, line width=0.5pt, -{Stealth[length=2.4pt]}] (0,0) -- (122,0) node[anchor=west, sxo lab] {$x$};
  \foreach \i in {0,1,2,3,4} {
    \fill[bkRed!80!black] ({10+25*\i},0) circle (1.6);
    \node[sxo lab, text=bkRed!80!black, anchor=north] at ({10+25*\i},-2.4) {$E_{\i}$};}
  \foreach \i in {0,1,2,3} {
    \filldraw[fill=white, draw=bkBlue!80!black, line width=0.9pt] ({22.5+25*\i},-1.4) rectangle ({25.3+25*\i},1.4);}
  \node[sxo lab, text=bkBlue!80!black, anchor=south] at (23.9,2.6) {$h_{\frac12}$};
  \node[sxo lab, text=bkBlue!80!black, anchor=south] at (48.9,2.6) {$h_{\frac32}$};
  \node[sxo lab, text=bkBlue!80!black, anchor=south] at (73.9,2.6) {$h_{\frac52}$};
  \node[sxo lab, text=bkBlue!80!black, anchor=south] at (98.9,2.6) {$h_{\frac72}$};
  \draw[sxo dim] (10,-9) -- (35,-9); \node[sxo lab, anchor=north] at (22.5,-10) {$\Delta x$};
  \node[sxo lab, anchor=west, align=left] at (0,14) {●：電場を置く点\qquad □：磁場を置く点（半マスずれた所）};
\end{tikzpicture}
\bkfigcaption{1次元の FDTD 法のマス。電場（●）と磁場（□）を、半マスずらして交互に置く。時間も、電場の時刻と磁場の時刻を半歩ずらして、かわるがわる進める。}
\label{ph:電-14:fig-yee}
\end{bkfigure}

\begin{phcode}[Python]
N = 10
E = [0]*N          \# 電場 E[0]〜E[9]
h = [0]*N          \# h[i] は E[i] と E[i+1] の間の点
E[4], E[5], E[6] = 2, 4, 2
def show(n):
    print(n, ' '.join('\%2d' \% v for v in E), '|',
             ' '.join('\%2d' \% v for v in h))
show(0)
for n in range(1, 4):
    for i in range(N-1):
        h[i] = h[i] - (E[i+1] - E[i])
    for i in range(1, N):
        E[i] = E[i] - (h[i] - h[i-1])
    show(n)
\end{phcode}
\begin{phcodeout}
0  0  0  0  0  2  4  2  0  0  0 |  0  0  0  0  0  0  0  0  0  0
1  0  0  0  2  2  0  2  2  0  0 |  0  0  0 -2 -2  2  2  0  0  0
2  0  0  2  2  0  0  0  2  2  0 |  0  0 -2 -2  0  0  2  2  0  0
3  0  2  2  0  0  0  0  0  2  2 |  0 -2 -2  0  0  0  0  2  2  0
\end{phcodeout}
\noindent 左はしの数が歩、｜の左が電場 $E$、右が $h$（$\eta_0\times$磁場の強さ）です。

はじめに電場の山（2、4、2）だけを置くと、山は高さ2の2つの山（2、2）に分かれて、1歩に1マスずつ左右へ進みます。8-13⑧ の弦の表とそっくりです。右へ進む山では $h$ と $E$ が同じ符号、左へ進む山では逆の符号です。$\vec S=\vec E\times\vec H$ の向きが、ちょうど波の進む向きになっているのです。1歩で $\Delta x$ 進み、1歩の時間は $\Delta t=\Delta x/c$ なので、山は光の速さで進みます。はじめに $h$ を $E$ と同じ形に置いておくと、山は右へだけ進みます。

\begin{phyomi}[$c\,\Delta t\leqq\Delta x$（CFL 条件）]
8-13 の弦の表計算と同じく、1歩のあいだに波がマス1つより遠くへ進むと、計算はこわれます。光は1秒に \qty{3.00e8}{m} 進むので、\qty{1}{cm} のマスでは、1歩を $3.3\times10^{-11}$ 秒（33 ピコ秒）以下にしなければなりません（例題12）。電子レンジの中（数十 cm）を光が1往復するだけで、何十歩も要ります。3次元で、もっと細かいマスで計算するには、大きなコンピュータが要るのです。
\end{phyomi}

%% ---------------- たいせつ・ことば・記号 ----------------
\begin{taisetsu}
\begin{itemize}[leftmargin=1.2\zw, itemsep=2pt, topsep=0pt, label={\color{bkNavy}●}]\raggedright
\item 変位電流 $I_{\mathrm d}=\varepsilon_0\,\dd\Phi_E/\dd t$。変わる電場は、電流と同じように磁場の渦を作る。充電中のコンデンサーでは、極板の間の変位電流は導線の電流と等しい。
\item マクスウェル方程式：電荷は電場のわき出し、磁場にわき出しはない、変わる磁場は電場の渦を、電流と変わる電場は磁場の渦を作る。時間で変わらなければ、電場と磁場は別々の式に分かれる。
\item 物の中では $\vec D$（自由電荷）と $\vec H$（自由電流）で書く。ふつうの物では $\varepsilon_0\to\varepsilon$、$\mu_0\to\mu$。
\item 電磁場のエネルギー密度 $u=\tfrac12\varepsilon_0E^2+B^2/(2\mu_0)$。エネルギーの流れはポインティングベクトル $\vec S=\vec E\times\vec H$〔W/m²〕。抵抗線の表面から入る量を合わせると $VI$。
\item 電場と磁場は、見る人の動きで入れかわる（$\vec E'=\vec E+\vec v\times\vec B$）。
\item FDTD 法：電場と磁場を半マスずらして置き、かわるがわる進める。$c\,\Delta t\leqq\Delta x$。
\end{itemize}
\end{taisetsu}

\begin{kotoba}\small
\bkword{変位電流}{電場の流束が変わる速さに $\varepsilon_0$ をかけたもの。電流と同じように磁場を作る。}
\bkword{アンペール・マクスウェルの法則}{アンペールの法則に、変位電流を足した式。}
\bkword{マクスウェル方程式}{電場と磁場のわき出しと渦を決める4つの式。}
\bkword{自由電荷・自由電流}{物の分極や磁化によらない、置いた電荷と導線を流れる電流。}
\bkword{エネルギー密度}{\qty{1}{m^3} あたりにたくわえられたエネルギー。}
\bkword{ポインティングベクトル}{電磁場のエネルギーの流れの向きと強さを表す矢印 $\vec E\times\vec H$。}
\bkword{電磁場}{電場と磁場を合わせた1つの場。見る人によって、電場と磁場の分かれ方が変わる。}
\bkword{磁気単極子}{N極だけ、またはS極だけをもつ粒。まだ見つかっていない。}
\bkword{FDTD 法}{マクスウェル方程式を、空間と時間のマスの差で解く方法。}
\end{kotoba}
\bkkigocards{\bkkigocard{$\vec S$}{エス}{ポインティングベクトル〔W/m²〕。電磁場のエネルギーの流れ。電-14・15 では面積を $A$ と書く}{Poynting vector}{電-14}}{\bkkigocard{$I_{\mathrm d}$}{アイ・ディー}{変位電流〔A〕。$\varepsilon_0\,\dd\Phi_E/\dd t$}{displacement current}{電-14}}

\begin{senmonsho}[マクスウェル方程式とポインティングの定理]
\begin{bkhonbun}
媒質中のマクスウェル方程式は $\nabla\cdot\symbfit{D}=\rho$，$\nabla\cdot\symbfit{B}=0$，$\nabla\times\symbfit{E}=-\partial\symbfit{B}/\partial t$，$\nabla\times\symbfit{H}=\symbfit{J}+\partial\symbfit{D}/\partial t$ である。第4式右辺第2項 $\partial\symbfit{D}/\partial t$ を変位電流密度という。これより $-\partial w/\partial t=\nabla\cdot\symbfit{S}+\symbfit{J}\cdot\symbfit{E}$（ポインティングの定理）が導かれる。ここで $w$ は電磁界のエネルギー密度，$\symbfit{S}=\symbfit{E}\times\symbfit{H}$ はポインティングベクトルである。
\end{bkhonbun}
\bkyaku{太字の $\symbfit{D}$・$\symbfit{E}$ などは、この本の矢印の $\vec D$・$\vec E$。$\nabla\cdot$ は div、$\nabla\times$ は rot（8-14⑪）。$\rho$ は自由電荷の密度、$\symbfit{J}$ は電流密度（この本の $\vec j$）。$w$ はこの本の $u$。「電磁界」は、電気工学の本の言い方で、この本の「電磁場」と同じもの（電界＝電場、磁界は $\vec H$ を指すことが多い）。最後の式は、⑦の保存の式を、場所ごとに書いたもの。$\symbfit{J}\cdot\symbfit{E}$ が、\qty{1}{m^3} あたりのジュール熱。}
\end{senmonsho}

%% ---------------- 例題 ----------------
\bksection{reidai}{例題でたしかめよう}

\noindent 定数は $\varepsilon_0=8.85\times10^{-12}$\,F/m、$\mu_0=4\pi\times10^{-7}$\,N/A²、$c=3.00\times10^8$\,m/s とします。$\mu_0/(2\pi)=2\times10^{-7}$\,N/A² です。

\begin{phreidai}[充電中のコンデンサーの、極板の間の磁場]
\label{ph:電-14:ex1}半径 \qty{5.0}{cm} の円い2枚の極板でできたコンデンサーを、\qty{1.0}{A} の電流で充電しています。(1) 極板の間を流れる変位電流を求めましょう。(2) 極板の縁（中心軸から \qty{5.0}{cm}）と、中心軸から \qty{2.5}{cm} の所の、極板の間の磁場の大きさを求めましょう。
\phbamen
②の図\ref{ph:電-14:fig-futatsu}の場面です。極板の間には電荷は流れていませんが、電場が強くなっていきます。
\phmushi
極板の縁の外へはみ出す電場（電-3 ④）を無視し、極板の間の電場は一様とします。導線は中心軸にそってまっすぐとします。
\phzu
極板の半径 $R=\qty{0.050}{m}$、充電の電流 $I=\qty{1.0}{A}$、中心軸からの距離 $r$、変位電流 $I_{\mathrm d}$、極板の間の磁場 $B$。
\phriyu
変位電流は③の $I_{\mathrm d}=\varepsilon_0\,\dd\Phi_E/\dd t$。磁場は、中心軸の周りの円い道（対称なので、道の上で $B$ は同じ大きさ）にアンペール・マクスウェルの法則を使う。
\phtateru
\[ I_{\mathrm d}=\varepsilon_0\,\frac{\dd\Phi_E}{\dd t}=\frac{\dd Q}{\dd t}=I,\qquad 2\pi r\times B=\mu_0\times I_{\mathrm d}\times\frac{\pi r^2}{\pi R^2} \]
（道が囲む変位電流は、全体の $(r/R)^2$。電場が一様なので、流束は面積に比例する。）
\phtoku
(1) $I_{\mathrm d}=\qty{1.0}{A}$。(2) $B=\dfrac{\mu_0I_{\mathrm d}}{2\pi}\times\dfrac{r}{R^2}$。縁（$r=R$）では
\[ B=2\times10^{-7}\times\frac{1.0}{0.050}=4.0\times10^{-6}\ \mathrm{T} \]
$r=\qty{0.025}{m}$ では $2\times10^{-7}\times1.0\times0.025\div0.050^2=2.0\times10^{-6}$\,T。
\phtashikame
単位：(N/A²)×A÷m＝N/(A·m)＝T。大きさ：地球の磁場（約 $4.6\times10^{-5}$\,T：電-6 ⑫）の十分の1ほどで、小さいが、はかれる大きさ。縁の値は、\qty{1.0}{A} の導線から \qty{5.0}{cm} の所と同じ（電-6 ⑤）。極端な場合：$r=0$ では0。
\kotae (1) \qty{1.0}{A}　(2) 縁 $4.0\times10^{-6}$\,T、\qty{2.5}{cm} で $2.0\times10^{-6}$\,T
\end{phreidai}
\begin{jibun}
半径 \qty{4.0}{cm} の円い極板のコンデンサーを \qty{0.50}{A} で充電しています。極板の縁と、中心軸から \qty{2.0}{cm} の所の、極板の間の磁場の大きさを求めましょう。
\bkjans{縁：$2\times10^{-7}\times0.50\div0.040=2.5\times10^{-6}$\,T。\qty{2.0}{cm}：$2\times10^{-7}\times0.50\times0.020\div0.040^2$ $\approx1.3\times10^{-6}$\,T}
\end{jibun}

\begin{phreidai}[電圧を上げる速さと、変位電流]
\label{ph:電-14:ex2}\qty{10}{cm}×\qty{10}{cm} の2枚の極板を、\qty{1.0}{mm} はなした平行板コンデンサーがあります。極板の間の電圧を、1秒に \qty{1000}{V} の割合で上げていきます。導線を流れる電流と、極板の間の変位電流を、それぞれの式で求め、等しいことを確かめましょう。
\phbamen
③の、電流と変位電流のつながりを、数で確かめます。
\phmushi
極板の縁で電場がはみ出すことを無視し、極板の間は真空（空気）とします。
\phzu
面積 $A=0.10\times0.10=\qty{0.010}{m^2}$、すき間 $d=\qty{1.0e-3}{m}$、電圧の上がる速さ $\dd V/\dd t=\qty{1000}{V/s}$。
\phriyu
導線の電流は $Q=CV$ を時間で微分した $I=C\,\dd V/\dd t$（電-3 ④）。変位電流は、$E=V/d$ から $\Phi_E=EA$ を作り、$\varepsilon_0\,\dd\Phi_E/\dd t$。
\phtateru
\[ C=\frac{\varepsilon_0A}{d},\qquad I=C\,\frac{\dd V}{\dd t},\qquad I_{\mathrm d}=\varepsilon_0A\,\frac{\dd E}{\dd t}=\varepsilon_0A\times\frac{1}{d}\,\frac{\dd V}{\dd t} \]
\phtoku
\[ C=\frac{8.85\times10^{-12}\times0.010}{0.0010}=8.85\times10^{-11}\ \mathrm{F}\ (\qty{88.5}{pF}) \]
\[ I=8.85\times10^{-11}\times1000=8.85\times10^{-8}\ \mathrm{A} \]
電場の増える速さは $1000\div0.0010=1.0\times10^6$\,V/(m·s) なので
\[ I_{\mathrm d}=8.85\times10^{-12}\times0.010\times1.0\times10^6=8.85\times10^{-8}\ \mathrm{A} \]
\phtashikame
単位：F×(V/s)＝C/s＝A。2つの式は、$C=\varepsilon_0A/d$ を入れると同じ式になる。大きさ：電場は1秒に $10^6$\,V/m（空気の絶縁が破れる強さの3分の1ほど）も増えるのに、変位電流は $10^{-7}$\,A ほどで、とても小さい。極端な場合：電圧が変わらなければ、どちらも0。
\kotae どちらも約 $8.9\times10^{-8}$\,A（\qty{89}{nA}）
\end{phreidai}
\begin{jibun}
\qty{20}{cm}×\qty{20}{cm} の極板を \qty{2.0}{mm} はなしたコンデンサーの電圧を、1秒に \qty{500}{V} の割合で上げます。導線の電流（＝変位電流）を求めましょう。
\bkjans{$C=8.85\times10^{-12}\times0.040\div0.0020$ $=1.77\times10^{-10}$\,F。$I=1.77\times10^{-10}\times500$ $\approx8.9\times10^{-8}$\,A}
\end{jibun}

\begin{phreidai}[海の中で、電波が届かないわけ]
\label{ph:電-14:ex3}海水の中で、周波数 \qty{1}{MHz} の電場がゆれています。海水は電気を通すので、ふつうの電流（伝導電流）も流れます。変位電流と伝導電流の大きさの比を求めましょう。海水の導電率を \qty{4.3}{S/m}、比誘電率を80とします。
\phbamen
潜水艦と陸の間の通信や、水の中のドローンの通信を考えます（電-15 ⑩）。
\phmushi
海水の性質が周波数によらないとし、場所による違いを無視します。電場は正弦波 $E=E_0\sin\omega t$ でゆれるとします（$t$ は時刻、$\omega=2\pi f$）。
\phzu
導電率 $\sigma=\qty{4.3}{S/m}$、誘電率 $\varepsilon=80\varepsilon_0$、周波数 $f=\qty{1.0e6}{Hz}$。\qty{1}{m^2} あたりの伝導電流 $j_{\text{伝}}$ と変位電流 $j_{\mathrm d}$。
\phriyu
伝導電流は $j=\sigma E$（電-4 ③）。物の中の変位電流は、⑥の $\partial D/\partial t$ で、$D=\varepsilon E$。$\sin\omega t$ を時間で微分すると $\omega\cos\omega t$ になるので（5-15）、ゆれの大きさは $\omega$ 倍。
\phtateru
\[ \frac{j_{\mathrm d}\ \text{の振幅}}{j_{\text{伝}}\ \text{の振幅}}=\frac{\varepsilon\omega E_0}{\sigma E_0}=\frac{\omega\varepsilon}{\sigma}=\frac{2\pi f\times80\varepsilon_0}{\sigma} \]
\phtoku
\[ \frac{2\pi\times1.0\times10^6\times80\times8.85\times10^{-12}}{4.3}\approx1.0\times10^{-3} \]
変位電流は、伝導電流の約1000分の1です。
\phtashikame
単位：(1/s)×(F/m)÷(S/m)＝F/(S·s)。F＝C/V、S＝A/V なので、C/(A·s)＝1（単位なし）。大きさ：海水の中では、電流のほとんどが伝導電流で、電場のエネルギーはすぐ熱になって消える。だから電波は海の中を遠くまで進めない。極端な場合：周波数を1000倍の \qty{1}{GHz} にすると比は約1。真水（$\sigma$ がずっと小さい）では、変位電流が主になる。
\kotae 約 $1.0\times10^{-3}$（変位電流は伝導電流の約1000分の1）
\end{phreidai}
\begin{jibun}
川の水（導電率 \qty{0.010}{S/m}、比誘電率80とする）の中で、\qty{1}{MHz} の電場がゆれるとき、変位電流と伝導電流の比を求めましょう。
\bkjans{$2\pi\times10^6\times80\times8.85\times10^{-12}\div0.010$ $\approx0.44$。変位電流は伝導電流の約0.44倍}
\end{jibun}

\begin{phreidai}[電場と磁場に、たくわえられるエネルギー]
\label{ph:電-14:ex4}(1) 電場 \qty{1.0e6}{V/m} の空間と、(2) 磁場 \qty{1.0}{T} の空間の、\qty{1}{m^3} あたりのエネルギーを求めましょう。(3) 磁場 \qty{1.0}{T} と同じエネルギー密度にするには、電場をいくらにしなければなりませんか。
\phbamen
\qty{1.0e6}{V/m} は、空気の絶縁が破れる手前（約 $3\times10^6$\,V/m）の強い電場、\qty{1.0}{T} は、モーターやスピーカーの磁石のすき間でよくある磁場です。
\phmushi
真空（空気）の中とし、物の分極・磁化を無視します。
\phzu
電場 $E=\qty{1.0e6}{V/m}$、磁場 $B=\qty{1.0}{T}$、エネルギー密度 $u_E$・$u_B$。
\phriyu
⑦のエネルギー密度の式の、電場の項と磁場の項を別々に計算する。
\phtateru
\[ u_E=\frac12\varepsilon_0E^2,\qquad u_B=\frac{B^2}{2\mu_0},\qquad \frac12\varepsilon_0E^2=u_B\ \ \text{から}\ \ E=\sqrt{\frac{2u_B}{\varepsilon_0}} \]
\phtoku
\[ u_E=\frac12\times8.85\times10^{-12}\times(1.0\times10^6)^2\approx4.4\ \mathrm{J/m^3} \]
\[ u_B=\frac{1.0^2}{2\times4\pi\times10^{-7}}\approx4.0\times10^5\ \mathrm{J/m^3} \]
\[ E=\sqrt{\frac{2\times4.0\times10^5}{8.85\times10^{-12}}}\approx3.0\times10^8\ \mathrm{V/m} \]
\phtashikame
単位：(F/m)×(V/m)²＝(C/V)×V²/m³＝J/m³。T²÷(N/A²)＝J/m³。大きさ：磁場のほうが約9万倍も多くのエネルギーをためられる。磁場と同じにするには、空気の絶縁の強さの100倍の電場が要り、無理。3の答えの $3.0\times10^8$ は、ちょうど $c\times B$（光の速さ×磁場）になっている（電-15 ④の $E=cB$ のとき、電場と磁場のエネルギーは等しい）。
\kotae (1) 約 \qty{4.4}{J/m^3}　(2) 約 $4.0\times10^5$\,J/m³　(3) 約 $3.0\times10^8$\,V/m
\end{phreidai}
\begin{jibun}
電場 \qty{3.0e6}{V/m} の空間と、磁場 \qty{0.50}{T} の空間の、\qty{1}{m^3} あたりのエネルギーを求めましょう。
\bkjans{電場：$\tfrac12\times8.85\times10^{-12}\times(3.0\times10^6)^2$ $\approx\qty{40}{J/m^3}$。磁場：$0.50^2\div(8\pi\times10^{-7})$ $\approx9.9\times10^4$\,J/m³}
\end{jibun}

\noindent モーターや変圧器が、電場ではなく磁場を使ってエネルギーを変えたり運んだりするのは、このためです（電-12）。

\begin{phreidai}[ドライヤーのヒーター線に入るエネルギー]
\label{ph:電-14:ex5}ドライヤーのヒーター線を、長さ \qty{0.50}{m}、半径 \qty{0.50}{mm} のまっすぐな線とします。両はしに \qty{100}{V} をかけると \qty{12}{A} 流れます。線の表面での電場・磁場の強さ・ポインティングベクトルの大きさを求め、表面全体から入るエネルギーが $VI$ に等しいことを確かめましょう。
\phbamen
⑧の抵抗線の場面です。ドライヤーの数（\qty{100}{V}・\qty{12}{A}・\qty{1200}{W}：電-4）を使います。
\phmushi
本物のヒーター線は巻いてありますが、まっすぐと考えます。線の太さは一様、周りのほかの導線の場を無視します。
\phzu
$V=\qty{100}{V}$、$I=\qty{12}{A}$、長さ $l=\qty{0.50}{m}$、半径 $a=\qty{5.0e-4}{m}$。表面の電場 $E$、磁場の強さ $H$、$S=EH$。
\phriyu
電場は線にそって一様で $E=V/l$。磁場の強さは、アンペールの法則で $H=I/(2\pi a)$。電場と磁場は直角なので、$S$ の大きさは $EH$。表面の面積は $2\pi al$。
\phtateru
\[ E=\frac{V}{l},\qquad H=\frac{I}{2\pi a},\qquad S=EH,\qquad P=S\times2\pi al \]
\phtoku
\[ E=\frac{100}{0.50}=200\ \mathrm{V/m},\qquad H=\frac{12}{2\pi\times5.0\times10^{-4}}\approx3.82\times10^3\ \mathrm{A/m} \]
\[ S=200\times3.82\times10^3\approx7.6\times10^5\ \mathrm{W/m^2} \]
表面の面積は $2\pi\times5.0\times10^{-4}\times0.50\approx1.57\times10^{-3}$\,m² なので
\[ P\approx7.64\times10^5\times1.57\times10^{-3}\approx1200\ \mathrm{W} \]
\phtashikame
単位：(V/m)×(A/m)＝W/m²。$VI=100\times12=\qty{1200}{W}$ と一致する。大きさ：$7.6\times10^5$\,W/m² は、大気の外で受ける太陽の光（\qty{1361}{W/m^2}：電-15 例題3）の500倍以上。細い線の表面に、それだけ強くエネルギーが流れこみ、線を赤く熱する。極端な場合：電流が0なら $H=0$ で、エネルギーは入らない。
\kotae $E=\qty{200}{V/m}$、$H\approx3.8\times10^3$\,A/m、$S\approx7.6\times10^5$\,W/m²。合計 \qty{1200}{W}（$=VI$）
\end{phreidai}
\begin{jibun}
豆電球のフィラメントを、長さ \qty{10}{cm}、半径 \qty{0.10}{mm} のまっすぐな線とします。\qty{1.5}{V} で \qty{0.30}{A} 流れるとき、表面のポインティングベクトルの大きさと、表面全体から入るエネルギーを求めましょう。
\bkjans{$E=1.5\div0.10=\qty{15}{V/m}$、$H=0.30\div(2\pi\times1.0\times10^{-4})$ $\approx478$\,A/m。$S\approx7.2\times10^3$\,W/m²。面積 $2\pi\times10^{-4}\times0.10$ をかけて \qty{0.45}{W}（$=1.5\times0.30$）}
\end{jibun}

\begin{phreidai}[同軸ケーブルの中を流れるエネルギー]
\label{ph:電-14:ex6}内側の導体の半径 \qty{0.45}{mm}、外側の導体の内側の半径 \qty{1.6}{mm} の同軸ケーブルで、\qty{10}{V} の直流を、\qty{50}{\ohm} の抵抗へ送ります。(1) 送る電力を求めましょう。(2) 中心から \qty{1.0}{mm} の所の、ポインティングベクトルの大きさを求めましょう。(3) 誘電体の中のポインティングベクトルをたし合わせると、(1) の電力になることを確かめましょう。
\phbamen
⑧の同軸ケーブルの場面です。寸法は、電-3 例題9・電-13 例題1 と同じ（特性インピーダンス約 \qty{50}{\ohm} のケーブル）です。
\phmushi
導体の抵抗を無視します（導体の中に電場はない）。誘電体の分極は電場の強さを変えますが、$V$ と $I$ を使って書くと、答えには入りません。
\phzu
$a=\qty{4.5e-4}{m}$、$b=\qty{1.6e-3}{m}$、$V=\qty{10}{V}$、$I=V/(\qty{50}{\ohm})$、中心からの距離 $r$。
\phriyu
誘電体の中の電場は、ガウスの法則から $1/r$ に比例し（電-3 ⑩）、内と外の導体の電位の差が $V$ になるように決まる：$E=V/(r\ln(b/a))$。磁場の強さは、アンペールの法則で $H=I/(2\pi r)$。2つは直角で、$\vec S$ はケーブルにそった向き。
\phtateru
\[ S=EH=\frac{V}{r\ln(b/a)}\times\frac{I}{2\pi r}=\frac{VI}{2\pi r^2\ln(b/a)} \]
半径 $r$、はば $\dd r$ の細い輪の面積は $2\pi r\,\dd r$ なので、全体は
\[ \int_a^b S\times2\pi r\,\dd r=\frac{VI}{\ln(b/a)}\int_a^b\frac{\dd r}{r}=\frac{VI}{\ln(b/a)}\times\ln\frac{b}{a} \]
\phtoku
(1) $I=10\div50=\qty{0.20}{A}$、$P=10\times0.20=\qty{2.0}{W}$。\par
(2) $\ln(1.6/0.45)\approx1.27$ なので
\[ S=\frac{10\times0.20}{2\pi\times(1.0\times10^{-3})^2\times1.27}\approx2.5\times10^5\ \mathrm{W/m^2} \]
(3) 積分の答えは $VI=\qty{2.0}{W}$ で、(1) と一致する（$\int\dd r/r=\ln r$：5-14）。
\phtashikame
単位：V×A÷(m²)＝W/m²。大きさ：内側の導線の近くほど強く流れ（$1/r^2$）、外の導体の近くでは約8分の1。誘電体の外（外の導体の外側）では $E=0$ で、エネルギーは流れない。極端な場合：$b\to a$（すき間が細い）でも、合計は $VI$ のまま。
\kotae (1) \qty{2.0}{W}　(2) 約 $2.5\times10^5$\,W/m²　(3) 合計 $VI=\qty{2.0}{W}$
\end{phreidai}
\begin{jibun}
同じケーブルで、\qty{5.0}{V} の直流を \qty{75}{\ohm} の抵抗へ送ります。送る電力と、中心から \qty{1.0}{mm} の所のポインティングベクトルの大きさを求めましょう。
\bkjans{$I=5.0\div75\approx0.067$\,A、$P\approx\qty{0.33}{W}$。$S=0.333\div(2\pi\times10^{-6}\times1.27)$ $\approx4.2\times10^4$\,W/m²}
\end{jibun}

\begin{phreidai}[コンデンサーに横から入るエネルギー]
\label{ph:電-14:ex7}半径 \qty{5.0}{cm} の円い極板を \qty{1.0}{mm} はなしたコンデンサーを、\qty{1.0}{A} で充電しています。電圧が \qty{100}{V} になった瞬間の、極板の縁でのポインティングベクトルの大きさを求め、縁の側面から入るエネルギーが、コンデンサーのエネルギー $\tfrac12CV^2$ の増える速さ $VI$ に等しいことを確かめましょう。
\phbamen
⑦の「身近な例では」の場面です（図\ref{ph:電-14:fig-condS}）。
\phmushi
極板の縁で電場がはみ出すことを無視します。極板の間は真空とします。
\phzu
$R=\qty{0.050}{m}$、$d=\qty{1.0e-3}{m}$、$V=\qty{100}{V}$、$I=\qty{1.0}{A}$。縁の電場 $E$、磁場の強さ $H$、縁の側面（高さ $d$ の円筒）の面積 $2\pi Rd$。
\phriyu
電場は $E=V/d$。縁の磁場は、例題1と同じく変位電流 $I$ 全体を囲むので $H=I/(2\pi R)$。$\vec S$ は縁から中心へ向かう。コンデンサーのエネルギー $\tfrac12CV^2$ の増える速さは、$I=C\,\dd V/\dd t$ を使って
\[ \frac{\dd}{\dd t}\Bigl(\frac12CV^2\Bigr)=CV\frac{\dd V}{\dd t}=VI \]
\phtateru
\[ S=\frac{V}{d}\times\frac{I}{2\pi R},\qquad P_{\text{入}}=S\times2\pi Rd=VI \]
\phtoku
\[ E=\frac{100}{0.0010}=1.0\times10^5\ \mathrm{V/m},\qquad H=\frac{1.0}{2\pi\times0.050}\approx3.18\ \mathrm{A/m} \]
$S=EH\approx3.2\times10^5$\,W/m²。
側面の面積 $2\pi\times0.050\times0.0010\approx3.14\times10^{-4}$\,m² をかけて $P_{\text{入}}\approx\qty{100}{W}$。$VI=100\times1.0=\qty{100}{W}$ と一致する。
\phtashikame
単位：W/m²×m²＝W。式の上でも、$2\pi R$ と $d$ が約分されて $VI$ になる。大きさ：このコンデンサーは $C\approx\qty{70}{pF}$ と小さく、\qty{1}{A} で充電すると1秒に $10^{10}$\,V 以上の速さで電圧が上がるので、\qty{100}{V} になるのは数ナノ秒（$10^{-9}$\,s）後。極端な場合：充電が終わると $I=0$ で $H=0$、エネルギーは流れこまない。
\kotae $S\approx3.2\times10^5$\,W/m²、側面から入るエネルギーは \qty{100}{W}（$=VI$）
\end{phreidai}
\begin{jibun}
例題7のコンデンサーを、\qty{0.50}{A} で充電していて、電圧が \qty{200}{V} になった瞬間に、縁の側面から入るエネルギーを求めましょう。
\bkjans{$E=200\div0.0010$ $=2.0\times10^5$\,V/m、$H=0.50\div(2\pi\times0.050)$ $\approx1.59$\,A/m、$S\approx3.2\times10^5$\,W/m²。側面の面積 $3.14\times10^{-4}$\,m² をかけて約 \qty{100}{W}（$=200\times0.50$）}
\end{jibun}

\begin{phreidai}[棒といっしょに動く人が見る電場]
\label{ph:電-14:ex8}磁場 \qty{0.10}{T} の中を、長さ \qty{0.50}{m} の金属の棒が、棒にも磁場にも直角な向きに \qty{2.0}{m/s} で動いています（電-8 の例題と同じ）。棒といっしょに動く人が見る、棒のある所の電場の大きさと、棒の両はしの電圧を求めましょう。
\phbamen
⑨の図\ref{ph:電-14:fig-futari}の場面です。
\phmushi
棒の太さ、棒の抵抗、棒の両はしにたまった電荷が作る電場を無視します（電流は流れていないとする）。速さは光の速さよりずっと小さい。
\phzu
$v=\qty{2.0}{m/s}$、$B=\qty{0.10}{T}$、$l=\qty{0.50}{m}$。動く人が見る電場 $E'$。
\phriyu
止まった人には電場がなく、動く人には $\vec E'=\vec v\times\vec B$ の電場が見える（⑨）。$\vec v$ と $\vec B$ は直角なので大きさは $vB$。電場が棒にそって一様なので、電圧は電場×長さ。
\phtateru
\[ E'=vB,\qquad V=E'l \]
\phtoku
$E'=2.0\times0.10=\qty{0.20}{V/m}$。$V=0.20\times0.50=\qty{0.10}{V}$。
\phtashikame
単位：(m/s)×T＝(m/s)×N/(A·m)＝N/C＝V/m。大きさ：電-8 ④ で、止まった人がローレンツ力から出した起電力 $vBl=\qty{0.10}{V}$ と同じ。見る人によって「電場」か「磁場の力」かは違っても、はかれる電圧は同じ。極端な場合：$v=0$ なら、どちらの人にも電場はない。
\kotae 電場 \qty{0.20}{V/m}、電圧 \qty{0.10}{V}
\end{phreidai}
\begin{jibun}
磁場 \qty{0.20}{T} の中を、長さ \qty{0.30}{m} の棒が、磁場と棒に直角な向きに \qty{5.0}{m/s} で動いています。棒といっしょに動く人が見る電場の大きさと、棒の両はしの電圧を求めましょう。
\bkjans{$E'=5.0\times0.20=\qty{1.0}{V/m}$、$V=1.0\times0.30=\qty{0.30}{V}$}
\end{jibun}

\begin{phreidai}[時間で変わらないとき、4つの式は何になる？]
\label{ph:電-14:ex9}マクスウェル方程式で、場が時間で変わらないとします。(1) 4つの式（積分形）はそれぞれどうなりますか。(2) (1) の式を使って、真空の中の \qty{1.0}{nC} の点電荷から \qty{0.10}{m} の所の電場と、まっすぐな \qty{10}{A} の電流から \qty{5.0}{cm} の所の磁場を求めましょう。
\phbamen
⑤の「時間で変わらないとき」を、数で確かめます（8-5⑥ と 8-10⑪ と同じ数）。
\phmushi
点電荷の大きさ、導線の太さを無視します。導線は十分長いとします。
\phzu
$Q=\qty{1.0e-9}{C}$、$r=\qty{0.10}{m}$、$I=\qty{10}{A}$、導線からの距離 $\qty{0.050}{m}$。
\phriyu
時間で変わらないので、$\dd\Phi_B/\dd t=0$、$\dd\Phi_E/\dd t=0$。（1）は電荷を中心とする球の袋、（4）は導線を囲む円の道を選ぶと、対称性から場の大きさが袋・道の上で同じになる（8-11⑦・⑬）。
\phtateru
(1)
\[ \oiint\vec E\cdot\vec n\,\dd A=\frac{Q_{\text{中}}}{\varepsilon_0},\qquad \oiint\vec B\cdot\vec n\,\dd A=0 \]
\[ \oint\vec E\cdot\dd\vec r=0,\qquad \oint\vec B\cdot\dd\vec r=\mu_0I_{\text{中}} \]
(2)
\[ E\times4\pi r^2=\frac{Q}{\varepsilon_0},\qquad B\times2\pi\times0.050=\mu_0I \]
\phtoku
$1/(4\pi\varepsilon_0)=8.99\times10^9$ を使うと
\[ E=\frac{Q}{4\pi\varepsilon_0r^2}=\frac{8.99\times10^9\times1.0\times10^{-9}}{0.10^2}\approx899\ \mathrm{N/C} \]
\[ B=\frac{2\times10^{-7}\times10}{0.050}=4.0\times10^{-5}\ \mathrm{T} \]
\phtashikame
単位：N/C＝V/m、T。大きさ：8-5⑥ の表の \qty{899}{N/C}、8-10⑪ の $4\times10^{-5}$\,T と同じ。時間で変わらなければ、電場は電荷だけ、磁場は電流だけで決まり、クーロンの法則（電-1）とアンペールの法則（電-6）にもどる。
\kotae (1) $\dd\Phi_B/\dd t$・$\dd\Phi_E/\dd t$ の項が消える（上の4つ）。(2) 電場 約 \qty{899}{N/C}、磁場 $4.0\times10^{-5}$\,T
\end{phreidai}
\begin{jibun}
例題9の(1)の式を使って、\qty{2.0}{nC} の点電荷から \qty{0.20}{m} の所の電場と、\qty{20}{A} の電流から \qty{10}{cm} の所の磁場を求めましょう。
\bkjans{$E=8.99\times10^9\times2.0\times10^{-9}\div0.20^2$ $\approx\qty{450}{N/C}$。$B=2\times10^{-7}\times20\div0.10=4.0\times10^{-5}$\,T}
\end{jibun}

\begin{phreidai}[水の中の電荷の周りの電場]
\label{ph:電-14:ex10}水（比誘電率80）の中に、\qty{1.0}{nC} の点電荷を置きました。\qty{0.10}{m} の所の電束密度 $D$ と電場 $E$ を求め、真空の中の電場と比べましょう。
\phbamen
⑥の物の中のガウスの法則を使います。
\phmushi
水の分子1つ1つの周りの細かな場を無視し、ならした場を考えます。水は動かず、電気を通さない（純粋な水）とします。電荷のすぐそばの水の性質の変化も無視します。
\phzu
$Q_{\text{自由}}=\qty{1.0e-9}{C}$、$r=\qty{0.10}{m}$、$\varepsilon_{\mathrm r}=80$。
\phriyu
⑥の（1）の $\vec D$ のガウスの法則は、自由電荷だけで書けるので、分極を数えなくてよい。球の袋を選ぶ。そのあと $D=\varepsilon_{\mathrm r}\varepsilon_0E$。
\phtateru
\[ D\times4\pi r^2=Q_{\text{自由}},\qquad E=\frac{D}{\varepsilon_{\mathrm r}\varepsilon_0} \]
\phtoku
\[ D=\frac{1.0\times10^{-9}}{4\pi\times0.10^2}\approx7.96\times10^{-9}\ \mathrm{C/m^2},\qquad E=\frac{7.96\times10^{-9}}{80\times8.85\times10^{-12}}\approx11\ \mathrm{V/m} \]
真空なら $7.96\times10^{-9}\div(8.85\times10^{-12})\approx\qty{899}{V/m}$ で、水の中は80分の1。
\phtashikame
単位：C/m² ÷ (F/m)＝(C/m²)×(V·m/C)＝V/m。大きさ：水の分子が電荷のほうへ向きをそろえ（分極）、電荷の周りに逆の符号の電荷が集まって、電場を弱める。食塩が水に溶けやすいのは、イオンどうしの引き合う力が水の中で弱まるからでもある。極端な場合：$\varepsilon_{\mathrm r}=1$ なら真空と同じ。
\kotae $D\approx8.0\times10^{-9}$\,C/m²、$E\approx\qty{11}{V/m}$（真空の約 \qty{899}{V/m} の80分の1）
\end{phreidai}
\begin{jibun}
例題10の電荷を、比誘電率2.3の油の中に置きました。\qty{0.10}{m} の所の電場を求めましょう。
\bkjans{$E=7.96\times10^{-9}\div(2.3\times8.85\times10^{-12})$ $\approx\qty{390}{V/m}$}
\end{jibun}

\begin{phreidai}[2つの面で、同じ磁場になることを確かめる]
\label{ph:電-14:ex11}例題1のコンデンサー（極板の半径 \qty{5.0}{cm}、\qty{1.0}{A} で充電）で、中心軸から $r=\qty{10}{cm}$ の円い道 $C$ を、極板の間の高さにとります。(1) 道 $C$ に張った平らな円板を面にして、道の上の磁場を求めましょう。(2) 道 $C$ から導線のほうへふくらませ、導線につらぬかれる袋の形の面にして、同じ磁場が出ることを確かめましょう。(3) 中心軸からの距離 $r$ と磁場 $B$ のグラフの形を説明しましょう。
\phbamen
②③の「面の選び方」を、極板の外の道で確かめます（図\ref{ph:電-14:fig-Br}）。
\phmushi
極板の縁で電場がはみ出すことを無視します。導線は中心軸にそってまっすぐとします。
\phzu
$R=\qty{0.050}{m}$、$I=\qty{1.0}{A}$、$r=\qty{0.10}{m}$。
\phriyu
アンペール・マクスウェルの法則では、面を通りぬける「電流＋変位電流」で数える。④で見たように、ふちが同じなら、どの面でも同じ値になるはず。
\phtateru
(1) 平らな円板（半径 \qty{10}{cm}）は極板の間を通り、電場の流束全部（変位電流 $I$）につらぬかれる。電流は0。(2) 袋の面は導線につらぬかれ（電流 $I$）、極板の間は通らない（変位電流0）。どちらも
\[ 2\pi r\times B=\mu_0\times I \]
\phtoku
$B=2\times10^{-7}\times1.0\div0.10=2.0\times10^{-6}$\,T。(3) $r\leqq R$ では $B=\mu_0Ir/(2\pi R^2)$ で $r$ に比例（例題1）、$r\geqq R$ では $\mu_0I/(2\pi r)$ で $r$ に反比例。縁でいちばん大きい山の形（図\ref{ph:電-14:fig-Br}）。
\phtashikame
単位：T。大きさ：(1)(2) とも同じ $2.0\times10^{-6}$\,T で、つじつまが合う。極板の外では、導線の周りの磁場（電-6 ⑤）と区別がつかない。極端な場合：$r$ が大きいほど0に近づく。
\kotae (1)(2) どちらも $2.0\times10^{-6}$\,T　(3) 縁までは $r$ に比例、縁の外では $r$ に反比例
\end{phreidai}
\begin{jibun}
\qty{2.0}{A} で充電しているコンデンサーの導線から \qty{3.0}{cm} の所の磁場を、導線につらぬかれる面と、極板の間を通る面の、2通りの考え方で求めましょう。
\bkjans{どちらも $2\pi\times0.030\times B=\mu_0\times2.0$ で、$B=2\times10^{-7}\times2.0\div0.030$ $\approx1.3\times10^{-5}$\,T}
\end{jibun}

\begin{phreidai}[FDTD のマスと、時間の1歩]
\label{ph:電-14:ex12}1次元の FDTD 法で、マスの大きさを $\Delta x=\qty{1.0}{cm}$ にとります。(1) 計算がこわれない時間の1歩 $\Delta t$ のいちばん大きい値を求めましょう。(2) $c\,\Delta t=\Delta x$ にとって100歩進めると、パルスは何 m 進みますか。そのとき、何秒たったことになりますか。
\phbamen
⑩の計算の準備です。
\phmushi
真空の中の1次元の波とし、はしでの反射を考えません。
\phzu
$\Delta x=\qty{0.010}{m}$、$c=\qty{3.00e8}{m/s}$、歩数100。
\phriyu
CFL 条件 $c\,\Delta t\leqq\Delta x$（⑩・8-13⑧）。$c\,\Delta t=\Delta x$ のとき、パルスは1歩に1マス進む（⑩の表）。
\phtateru
\[ \Delta t\leqq\frac{\Delta x}{c},\qquad \text{進む長さ}=100\times\Delta x,\qquad \text{時間}=100\times\Delta t \]
\phtoku
(1) $\Delta t\leqq0.010\div(3.00\times10^8)\approx3.3\times10^{-11}$\,s。(2) $100\times0.010=\qty{1.0}{m}$。時間は $100\times3.33\times10^{-11}\approx3.3\times10^{-9}$\,s。
\phtashikame
単位：m÷(m/s)＝s。大きさ：光は \qty{1}{m} を約 3.3 ナノ秒で進むので合う。極端な場合：マスを半分にすると、1歩も半分になり、同じ長さを進むのに4倍（1次元では2倍のマスと2倍の歩数）の計算が要る。
\kotae (1) $3.3\times10^{-11}$\,s　(2) \qty{1.0}{m}、$3.3\times10^{-9}$\,s
\end{phreidai}
\begin{jibun}
マスの大きさを \qty{2.0}{mm} にしたとき、時間の1歩のいちばん大きい値と、$c\,\Delta t=\Delta x$ で300歩進めたときにパルスが進む長さを求めましょう。
\bkjans{$\Delta t\leqq2.0\times10^{-3}\div(3.00\times10^8)$ $\approx6.7\times10^{-12}$\,s。$300\times\qty{2.0}{mm}=\qty{0.60}{m}$}
\end{jibun}


%% ---------------- まちがえやすいところ ----------------
\bksection{tsumazuki}{まちがえやすいところをチェック}
\begin{tsumazuki}{変位電流を、電荷が極板の間を流れることと考える}
\bkng{充電中は、極板の間を電荷が飛んでいるので、変位電流が流れる。}
\bkok{極板の間を通る電荷はない。電場が強くなっていくこと（$\dd\Phi_E/\dd t$）を、電流のなかまとして数えたのが変位電流。}
変位電流は、名前に「電流」とありますが、電荷の流れではありません。真空の中でも、電場が変われば変位電流があり、磁場の渦を作ります。
\end{tsumazuki}
\begin{tsumazuki}{充電が終わったあとも、磁場があると考える}
\bkng{コンデンサーには電場があるので、極板の間に磁場ができ続ける。}
\bkok{磁場の渦を作るのは、電場の「変化」。充電が終わって電場が変わらなくなれば、変位電流は0で、磁場もなくなる。}
変位電流は $\Phi_E$ そのものではなく、その変わる速さ $\dd\Phi_E/\dd t$ に比例します。同じように、（3）の電場の渦を作るのも、磁場そのものではなく磁場の変化です。
\end{tsumazuki}
\begin{tsumazuki}{エネルギーは導線の中を流れると考える}
\bkng{ヒーター線のエネルギーは、電流といっしょに導線の中を流れてくる。}
\bkok{エネルギーは、ポインティングベクトル $\vec E\times\vec H$ の向きに、導線の周りの空間を流れ、ヒーター線の表面から中へ入る。}
導線の中の電子の速さは、1秒に1 mm より小さいほどゆっくりです（電-4 ②）。それでも、スイッチを入れたとたんに電球がつくのは、エネルギーが電子といっしょに運ばれるのではなく、導線の周りの電磁場を通って、光の速さに近い速さで伝わるからです。
\end{tsumazuki}
\begin{tsumazuki}{$\vec S$ の向きを、外積の順番を逆にして決める}
\bkng{$\vec H$ から $\vec E$ へ右ねじを回して、$\vec S$ の向きを決める。}
\bkok{$\vec S=\vec E\times\vec H$ なので、$\vec E$ から $\vec H$ へ回す。}
外積は、順番を入れかえると向きが逆になります（4-13）。抵抗線で確かめると、$\vec E$ から $\vec H$ へ回すと、$\vec S$ は線の中を向き、エネルギーが入ってきて熱になる、と正しく読めます。逆にすると、抵抗線からエネルギーが出ていくことになり、おかしくなります。
\end{tsumazuki}

%% ---------------- 工学ではこう使う ----------------
\bksection{phgear}{工学ではこう使う}
\begin{phkougaku}[電磁場を計算して、機械を作る]
\phboxhead{① 電磁界解析}
スマホ、自動車のレーダー、電子レンジ、MRI（磁気共鳴画像）の装置などは、作る前にコンピュータでマクスウェル方程式を解いて、アンテナから出る電波や、中の電磁場の強さを確かめます。⑩の FDTD 法のほかに、形を小さな三角形や四面体に分けて解く有限要素法（力-16 と同じ考え）がよく使われます。

\phboxhead{② 電磁両立性（EMC）}
電気製品は、電磁場を外へもらして、ほかの製品の動きをじゃましてはいけません。また、外から来た電磁場で、自分がおかしくなってもいけません。これを\textbf{電磁両立性（EMC）}といいます。\bkidx{でんじりょうりつせい}{電磁両立性}回路の電流が変わると、⑤の（3）（4）で電磁場が生まれて外へ広がるので、配線の形や、金属のケースで囲む方法（電-15 ⑩）をくふうします。

\begin{phketa}[電磁場の数]
\phketaline{空気の絶縁が破れる電場}{約 $3\times10^6$\,V/m}
\phketaline{モーターのすき間の磁場}{\qty{1}{T} くらい}
\phketaline{充電中のコンデンサーの極板の間の磁場}{$10^{-6}$\,T くらい}
\phketaline{太陽の光のポインティングベクトル（大気の外）}{\qty{1361}{W/m^2}}
\phboxhead{答えが桁で合っているか}
エネルギー密度は、電場なら $\tfrac12\varepsilon_0E^2$ で、ふつうの電場では \qty{1}{m^3} あたり数 J 以下。磁場なら \qty{1}{T} で約 $4\times10^5$\,J/m³。ポインティングベクトルを面積でたし合わせた値が、回路の電力 $VI$ と合っているかで、計算を確かめられます。
\end{phketa}
\end{phkougaku}
\phstation{電磁気学,電磁界解析,EMC,高電圧工学}{マクスウェル方程式、変位電流、電束密度、磁界の強さ、ポインティングベクトル、電磁界のエネルギー、境界条件、FDTD 法、有限要素法}

\begin{column}{マクスウェルの「付け加えた項」}
1865年、イギリスのマクスウェルは、論文「電磁場の動力学的理論」で、それまでに見つかっていた電気と磁気の法則をまとめ、アンペールの法則に変位電流の項を付け加えました。そして、電場と磁場の変化が、1秒に約30万 km の速さで伝わる波になることを示しました。この速さが、当時はかられていた光の速さとほとんど同じだったことから、マクスウェルは「光は電磁場の波である」と考えました。電磁波を実験で作って確かめたのは、ドイツのヘルツで、1888年ごろのことです（8-14⑤のコラム）。

変位電流の項は、実験で見つかったものではなく、つじつまを合わせるために考え出されたものでした。その項から、まだだれも見たことのない「電磁波」が予言され、20年あまりのちに本当に見つかったのです。このあとの電-15 では、マクスウェルの式から電磁波を導きます。
\end{column}

\begin{column}{3次元のマス：イーの箱}
\noindent\begin{minipage}[c]{0.56\linewidth}\raggedright
⑩では1次元で、電場と磁場を半マスずらして置きました。3次元では、小さな立方体（イーの箱）の\textbf{辺のまん中}に電場の成分を、\textbf{面のまん中}に磁場の成分を置きます。こうすると、1つの面の周りの4本の辺の電場をたすと、（3）のファラデーの法則の「道を1周する電場」になり、その面のまん中の磁場を進められます。同じように、磁場の4つの点を1周たすと（4）になり、電場を進められます。積分形の4つの式を、小さな箱でそのまま計算しているのです。アンテナの設計などでは、この箱をとてもたくさん並べて、コンピュータで計算します。
\end{minipage}\hfill
\begin{minipage}[c]{0.4\linewidth}\centering
\begin{tikzpicture}[x=1mm, y=1mm]
  \draw[bkInk, line width=0.6pt] (0,0) rectangle (20,20);
  \draw[bkInk, line width=0.6pt] (20,0) -- (28,7) -- (28,27) -- (20,20) (0,20) -- (8,27) -- (28,27);
  \draw[bkInk!50, line width=0.5pt, dash pattern=on 1.5pt off 1pt] (0,0) -- (8,7) -- (28,7) (8,7) -- (8,27);
  % 辺のまん中の電場（実線の矢印）
  \draw[sxo E] (6,0) -- (14,0); \draw[sxo E] (20,6) -- (20,14); \draw[sxo E] (6,20) -- (14,20); \draw[sxo E] (0,6) -- (0,14);
  % 面のまん中の磁場（破線の矢印）
  \draw[sxo B] (10,10) -- (14.8,14.2);
  \draw[sxo B] (24,13.5) -- (30,13.5);
  \node[sxo lab, text=bkRed!80!black, anchor=north] at (10,-1.5) {辺に $\vec E$};
  \node[sxo lab, text=bkBlue!80!black, anchor=south] at (14,29) {面に $\vec H$};
  \draw[bkGray, line width=0.4pt] (14,28.8) -- (13,15);
\end{tikzpicture}
\end{minipage}
\end{column}

\begin{phsaki}{相対性理論と、場の理論へ}
この章では、電場と磁場が見る人で入れかわることを、速さが光より十分小さいときだけ見ました。速さが光に近いときの入れかわりは、特殊相対性理論で、電場と磁場を1つの「電磁場テンソル」にまとめて書きます（原-11）。電磁気学の専門の本では、マクスウェル方程式を、電位とベクトルポテンシャル（電-6 ⑦のもっと知りたい）の式に書きかえて、アンテナから出る電磁波を計算します（電-15 ⑧）。さらに進むと、電磁場を量子の考えで扱う「量子電磁力学」や、電磁気の力を、ほかの力（弱い力・強い力）と同じ形でまとめる「ゲージ理論」に進みます。
\end{phsaki}

%% ---------------- 章末問題 ----------------
\bkendproblems
\bkproblemgroup{基本}{3}
\bkproblem 半径 \qty{3.0}{cm} の円い極板のコンデンサーを \qty{0.30}{A} で充電しています。極板の間の変位電流と、極板の縁での磁場の大きさを求めましょう（\nobreak\bkreidairef{1}\nobreak）。
\bkpans{変位電流 \qty{0.30}{A}。$B=2\times10^{-7}\times0.30\div0.030=2.0\times10^{-6}$\,T}
\bkproblem 面積 \qty{0.020}{m^2}、すき間 \qty{0.50}{mm} の平行板コンデンサーの電圧を、1秒に \qty{200}{V} の割合で上げます。変位電流を求めましょう（\nobreak\bkreidairef{2}\nobreak）。
\bkpans{$C=8.85\times10^{-12}\times0.020\div(5.0\times10^{-4})$ $\approx3.5\times10^{-10}$\,F。$I_{\mathrm d}=3.54\times10^{-10}\times200$ $\approx7.1\times10^{-8}$\,A}
\bkproblem 電場 \qty{2.0e5}{V/m} と磁場 \qty{0.20}{T} が同じ所にあります。\qty{1}{m^3} あたりの電場のエネルギーと磁場のエネルギーを求めましょう（\nobreak\bkreidairef{4}\nobreak）。
\bkpans{電場：$\tfrac12\times8.85\times10^{-12}\times(2.0\times10^5)^2$ $\approx\qty{0.18}{J/m^3}$。磁場：$0.20^2\div(8\pi\times10^{-7})$ $\approx1.6\times10^4$\,J/m³}
\bkproblem \qty{0.80}{A} で充電しているコンデンサーの導線から \qty{2.0}{cm} の所の磁場の大きさを、導線につらぬかれる面と、極板の間を通る面の、2通りの考え方で求めましょう（\nobreak\bkreidairef{11}\nobreak）。
\bkpans{導線につらぬかれる面では電流 \qty{0.80}{A}、極板の間を通る面では変位電流 \qty{0.80}{A}。どちらも $B=2\times10^{-7}\times0.80\div0.020=8.0\times10^{-6}$\,T}
\bkproblem 半径 \qty{6.0}{cm} の円い極板のコンデンサーを \qty{1.5}{A} で充電しています。極板の間で、中心軸から \qty{3.0}{cm} の所と、極板の外の \qty{12}{cm} の所の磁場の大きさを求めましょう（\nobreak\bkreidairef{1}\nobreak・\nobreak\bkreidairef{11}\nobreak）。
\bkpans{\qty{3.0}{cm}：囲む変位電流は $1.5\times(3.0/6.0)^2=0.375$\,A で、$B=2\times10^{-7}\times0.375\div0.030=2.5\times10^{-6}$\,T。\qty{12}{cm}：$B=2\times10^{-7}\times1.5\div0.12=2.5\times10^{-6}$\,T}
\bkproblem 充電しているコンデンサーの＋の極板の電荷が、1秒に $2.0\times10^{-6}$\,C の割合で増えています。導線の電流、極板の間の電場の流束 $\Phi_E$ の増える速さ、変位電流を求めましょう（③・\nobreak\bkreidairef{2}\nobreak）。
\bkpans{電流 $2.0\times10^{-6}$\,A。$\dd\Phi_E/\dd t=I/\varepsilon_0=2.0\times10^{-6}\div(8.85\times10^{-12})$ $\approx2.3\times10^5$\,V·m/s。変位電流は $\varepsilon_0\times$ これ $=2.0\times10^{-6}$\,A（導線の電流と同じ）}
\bkproblem 1次元の FDTD 法で、マスの大きさを \qty{5.0}{mm} にとります。計算がこわれない時間の1歩のいちばん大きい値と、$c\,\Delta t=\Delta x$ で200歩進めたときにパルスが進む長さを求めましょう（\nobreak\bkreidairef{12}\nobreak）。
\bkpans{$\Delta t\leqq5.0\times10^{-3}\div(3.00\times10^8)$ $\approx1.7\times10^{-11}$\,s。$200\times\qty{5.0}{mm}=\qty{1.0}{m}$}
\bkproblemgroup{チャレンジ}{5}
\bkproblem 長さ \qty{1.0}{m}、半径 \qty{0.20}{mm} のニクロム線に \qty{12}{V} をかけると、\qty{2.0}{A} 流れました。表面のポインティングベクトルの大きさと、線に入るエネルギーを求めましょう（\nobreak\bkreidairef{5}\nobreak）。
\bkpans{$E=\qty{12}{V/m}$、$H=2.0\div(2\pi\times2.0\times10^{-4})$ $\approx1.6\times10^3$\,A/m。$S\approx1.9\times10^4$\,W/m²。面積 $2\pi\times2.0\times10^{-4}\times1.0$ をかけて \qty{24}{W}（$=12\times2.0$）}
\bkproblem 磁場 \qty{0.50}{T} の中を、長さ \qty{0.20}{m} の棒が \qty{3.0}{m/s} で、棒にも磁場にも直角に動いています。（1）棒といっしょに動く人が見る電場と、棒の両はしの電圧を求めましょう。（2）止まっている人は、この電圧を何の力で説明しますか（\nobreak\bkreidairef{8}\nobreak）。
\bkpans{（1）$E'=3.0\times0.50=\qty{1.5}{V/m}$、$V=1.5\times0.20=\qty{0.30}{V}$。（2）棒の中の電子が磁場から受けるローレンツ力 $q\vec v\times\vec B$。起電力は $vBl=\qty{0.30}{V}$ で同じ}
\bkproblem 例題6の同軸ケーブル（$a=\qty{0.45}{mm}$、$b=\qty{1.6}{mm}$）で \qty{10}{V}・\qty{0.20}{A} を送るとき、送るエネルギーの半分が通るのは、中心から何 mm より内側ですか（\nobreak\bkreidairef{6}\nobreak・5-14）。
\bkpans{半径 $r_{\frac12}$ までの合計は $VI\ln(r_{\frac12}/a)\div\ln(b/a)$。これが半分なので $\ln(r_{\frac12}/a)=\tfrac12\ln(b/a)$、$r_{\frac12}=\sqrt{ab}=\sqrt{0.45\times1.6}\approx\qty{0.85}{mm}$}
\bkproblem 1次元の FDTD 法で、⑩のプログラムの $E$ の山を（1、2、1）にして、3歩めまでの $E$ と $h$ を表計算か Python で計算し、山が2つに分かれて進むことを確かめましょう。分かれた山の高さはいくらですか。$h$ の符号から、それぞれの山が進む向きを説明しましょう（⑩）。
\bkpans{1歩めの $E$ は 0 0 0 1 1 0 1 1 0 0 で、（1、1）の山が2つに分かれ、1歩に1マスずつ左右へ進む。高さは1（もとのいちばん高い所の半分）。右へ進む山は $h$ が $E$ と同じ符号（＋1）、左へ進む山は逆の符号（－1）で、$\vec S=\vec E\times\vec H$ が進む向きを向く}

\begin{furikaeri}
\begin{checklist}
\item 充電しているコンデンサーの極板の間を流れる変位電流と、その周りの磁場を計算できる。\bkfurref{ph:電-14:sec3}
\item マクスウェル方程式を積分形と微分形で読み、電磁場のエネルギーの流れ（ポインティングベクトル）を計算できる。\bkfurref{ph:電-14:sec7}
\item 動く人から見ると電場と磁場が入れかわることを説明し、電磁場を数値で解く方法（FDTD）を表計算で動かせる。\bkfurref{ph:電-14:sec9}
\end{checklist}
\end{furikaeri}
\bkfukidashi{sora}{変わる電場が磁場を作り、変わる磁場が電場を作る。この2つがかわるがわる進むと、どうなるのかな。次の電-15 では、それが光と電波だとわかるよ。}
